\documentclass[12pt,a4paper]{article}
\usepackage[utf8]{inputenc}
\usepackage[T1]{fontenc}
\usepackage[main=english,french]{babel}
\usepackage{lmodern}
\usepackage{amsmath,amssymb,amsthm}
\usepackage{amsfonts}
\usepackage{mathrsfs}
\usepackage{geometry}
\usepackage{titlesec}
\usepackage{fancyhdr}
\usepackage{array}

\geometry{margin=1in}

\pagestyle{fancy}
\fancyhf{}
\fancyhead[C]{\thepage}
\renewcommand{\headrulewidth}{0pt}

\titleformat{\section}{\large\bfseries}{\thesection.}{0.5em}{}
\titleformat{\subsection}{\normalsize\bfseries}{\thesubsection.}{0.5em}{}

% Cayley's vertical bar (he uses \mid throughout for | in formulae such as K|, k|)
\providecommand{\Cmid}{\mid}

% French babel shorthands off so that : ! ? ; do not get treated specially
\providecommand{\noFrenchShorthands}{\shortsource packet{:!?;}}

\usepackage{graphicx}
\emergencystretch=3em
\tolerance=600
\begin{document}

\selectlanguage{english}

\begin{center}
{\LARGE\bfseries The Collected Mathematical Papers of}\\[0.5em]
{\LARGE\bfseries Arthur Cayley}\\[1em]
{\Large Volume I}\\[1em]
{\large PDF Pages 301--350 \quad (Book pages 279--326)}\\[2em]
\end{center}

\tableofcontents
\newpage

% ===========================================================================
% PDF page 301 -- Book page 279
% Continuation of paper [43] On the Differential Equations which occur in
% Dynamical Problems.
% ===========================================================================
\section*{[43] (continued)}
\section*{On the Differential Equations which occur in Dynamical Problems.}
\addcontentsline{toc}{section}{43. On the Differential Equations which occur in Dynamical Problems (continued)}

\textit{[From the Cambridge and Dublin Mathematical Journal, vol.~ii.~(1847), pp.~210--219.]}

\bigskip
\setcounter{page}{279}

Suppose that $u$ and $v$ continue to represent arbitrary functions of $x,y,z,\ldots$ but that the remaining functions $w,\ldots$ are such as to satisfy $\nabla = 0,\ldots$ (so that $w,\ldots$ may be considered as the constants introduced by obtaining all the integrals but one of the system of differential equations in $x,y,z,\ldots$), we have
\[
\frac{dMVU}{du} + \frac{dMVV}{dv} = V\left(\frac{dMX}{dx} + \frac{dMY}{dy} + \frac{dMZ}{dz} + \cdots\right).
\]

Also the only one of the transformed equations which remains to be integrated is
\[
du : dv = U : V, \quad\text{or}\quad V\,du - U\,dv = 0,
\]
(in which it is supposed that $U$ and $V$ are expressed by means of the other integrals in terms of $u$ and $v$).

Suppose $M$ can be so determined that
\[
\frac{dMX}{dx} + \frac{dMY}{dy} + \frac{dMZ}{dz} + \cdots = 0,
\]
($M$ is what Jacobi terms the multiplier of the proposed system of differential equations): then
\[
\frac{dMVU}{du} + \frac{dMVV}{dv} = 0,
\]
or $MV$ is the multiplier of $V\,du - U\,dv = 0$, so that
\[
\delta(MV(V\,du - U\,dv)) = \text{const.}
\]

Hence the theorem: ``Given a multiplier of the system of equations
\[
dx : dy : dz \ldots = X : Y : Z \ldots
\]
(the meaning of the term being defined as above), then if all the integrals but one of this system are known, the remaining integral depends upon a quadrature.''

Jacobi proceeds to discuss a variety of different systems of equations in which it is possible to determine the multiplier $M$. Among the most important of these may be considered the system corresponding to the general problem of Dynamics, which may be discussed under three different forms.

\subsection*{\S\,4. Lagrange's first form\footnotemark.}
\footnotetext{I have slightly modified the form so as to avoid the introduction of the masses, and to allow $x$ (for instance) to stand for any one of the coordinates of any of the points, instead of standing for a coordinate parallel to a particular axis.}

Let the whole series of coordinates, each of them multiplied by the square root of the corresponding mass, be represented by $x, y, \ldots$ and in the same way the whole series of forces, each of them multiplied by the square root of the corresponding mass, by $P, Q, \ldots$; then the equations of motion are
\[
\frac{d^2 x}{dt^2} = P + \lambda\frac{d\Theta}{dx} + \mu\frac{d\Phi}{dx} + \cdots,
\]
\[
\frac{d^2 y}{dt^2} = Q + \lambda\frac{d\Theta}{dy} + \mu\frac{d\Phi}{dy} + \cdots,
\]
where $\Theta = 0,\ \Phi = 0,\ldots$ are the equations of condition connecting the variables, and $\lambda, \mu,\ldots$ coefficients to be determined by substituting the values of $\frac{d^2x}{dt^2}$, \&c.\ in the equations $\frac{d^2\Theta}{dt^2}=0$, $\frac{d^2\Phi}{dt^2}=0$, \&c. It is supposed that as well $P, Q,\ldots$ as $\Theta, \Phi,\ldots$ are independent of the velocities.

In order to reduce these to an analogous form to that previously employed, we have only to write
\[
\frac{dx}{dt} = x',\quad \frac{dy}{dt} = y',\quad \ldots
\]
which gives
\[
dt : dx : dy : dz \ldots : dx' : dy' : dz' \ldots
\]
\[
= 1 : x' : y' : z' \ldots : X : Y : Z \ldots
\]

Supposing that $M$ is independent of $x', y', z',\ldots$ the equation on which it depends becomes immediately
\[
\delta M + M\left(\frac{dX}{dx'} + \frac{dY}{dy'} + \frac{dZ}{dz'} + \cdots\right) = 0,
\]
where for shortness
\[
\delta = \frac{d}{dt} + x'\frac{d}{dx} + y'\frac{d}{dy} + z'\frac{d}{dz} + \cdots.
\]

To reduce this we must first determine the values of $\lambda, \mu,\ldots,$ and for this we have
\[
\frac{d^2\Theta}{dt^2} = \frac{d^2\Theta}{dx^2}{x'}^2 + \cdots + \frac{d\Theta}{dx}\frac{dx'}{dt} + \cdots = 0,
\]
i.e.
\[
\frac{d\Theta}{dx}\frac{dx'}{dt} + \cdots + a{x'}^2 + b{y'}^2 + \cdots = 0,
\]
where for greater clearness an additional letter of the series $\Theta, \Phi,\ldots$ has been introduced, and where
\[
a = \left(\frac{d^2\Theta}{dx^2}\right),\quad h = \left(\frac{d^2\Theta}{dx\,dy}\right),\quad \ldots
\]

Hence differentiating with respect to $x'$,
\[
a\frac{d\Theta}{dx} + h\frac{d\Phi}{dx} + g\frac{d\Psi}{dx} + \cdots = 0,
\]
\[
h\frac{d\Theta}{dx} + b\frac{d\Phi}{dx} + f\frac{d\Psi}{dx} + \cdots = 0,
\]
\[
g\frac{d\Theta}{dx} + f\frac{d\Phi}{dx} + c\frac{d\Psi}{dx} + \cdots = 0;
\]
or representing by $K$ the determinant formed with the quantities $a, h, g, \ldots, h, b, f, \ldots, g, f, c, \ldots$ and by $A, H, G, \ldots, H, B, F, \ldots, G, F, C, \ldots$ the inverse system of coefficients, we have
\[
K\left(\lambda\frac{d\Theta}{dx} + \mu\frac{d\Phi}{dx} + \nu\frac{d\Psi}{dx} + \cdots\right) + K\frac{dx'}{dt} = 0,
\]
whence, multiplying by $\frac{d\Theta}{dx}, \frac{d\Phi}{dx}, \frac{d\Psi}{dx}, \ldots$ and adding,
\[
\lambda\left(\frac{d\Theta}{dx}\right)^2 + \mu\frac{d\Theta}{dx}\frac{d\Phi}{dx} + \cdots + \frac{dx'}{dt}\frac{d\Theta}{dx} = 0,
\]
and, forming the similar equations with the remaining variables and adding,
\[
\Sigma\Bigl\{ A\delta a + B\delta b + C\delta c + \cdots + 2F\delta f + 2G\delta g + 2H\delta h + \cdots + K\Bigl(\frac{dx'}{dx} + \frac{dy'}{dy} + \frac{dz'}{dz}+\cdots\Bigr)\Bigr\} = 0;
\]
i.e.
\[
\frac{dX}{dx'} + \frac{dY}{dy'} + \frac{dZ}{dz'} + \cdots = 0.
\]

Thus the equation in $M$ reduces itself to
\[
K\delta M - M\delta K = 0,
\]
which is satisfied by $M = K$. It may be remarked that $K$ reduces itself to the sum of the squares of the different functional determinants formed with the differential coefficients of $\Theta, \Phi,\ldots$ with respect to the different combinations of as many variables out of the series $x, y, \ldots$.

\subsection*{\S\,5. Lagrange's second form.}

Here the equations of motion are assumed to be
\[
\frac{d}{dt}\frac{dT}{dx'} - \frac{dT}{dx} - P = 0,
\]
\[
\frac{d}{dt}\frac{dT}{dy'} - \frac{dT}{dy} - Q = 0,
\]
\[
\frac{d}{dt}\frac{dT}{dz'} - \frac{dT}{dz} - R = 0,
\]
where $2T$ represents the \emph{vis viva} of the system, $x, y, z,\ldots$ are the independent variables on which the solution of the problem depends, and $x', y', z',\ldots$ their differential coefficients with respect to the time. It is assumed as before $P, Q, R\ldots$ do not contain $x', y', z',\ldots$.

Suppose these equations give
\[
dt : dx : dy : dz \ldots : dx' : dy' : dz' \ldots
\]
\[
= 1 : x' : y' : z' \ldots : X : Y : Z \ldots ;
\]
then the equation which determines the multiplier $M$ takes as before the form
\[
\delta M + M\left(\frac{dX}{dx'} + \frac{dY}{dy'} + \frac{dZ}{dz'} + \cdots\right) = 0.
\]

To reduce this equation, substituting for $T$ its value which is of the form
\[
T = \tfrac12\bigl(a{x'}^2 + b{y'}^2 + c{z'}^2 \ldots + 2f y' z' + 2g z' x' + 2h x' y' \ldots \bigr),
\]
and putting for shortness
\[
L = a x' + h y' + g z' \ldots,
\]
\[
M = h x' + b y' + f z' \ldots,
\]
\[
N = g x' + f y' + c z' \ldots,
\]
the equations which determine $X, Y, Z, \ldots$ are
\[
aX + hY + gZ \ldots + \delta L - \frac{dT}{dx} - P = 0,
\]
\[
hX + bY + fZ \ldots + \delta M - \frac{dT}{dy} - Q = 0,
\]
\[
gX + fY + cZ \ldots + \delta N - \frac{dT}{dz} - R = 0.
\]

Hence, differentiating with respect to $x'$,
\[
a\frac{dX}{dx'} + h\frac{dY}{dx'} + g\frac{dZ}{dx'} \ldots + \delta a = 0,
\]
\[
h\frac{dX}{dx'} + b\frac{dY}{dx'} + f\frac{dZ}{dx'} \ldots + \delta h + a' = 0,
\]
\[
g\frac{dX}{dx'} + f\frac{dY}{dx'} + c\frac{dZ}{dx'} \ldots + \delta g = 0,
\]
or representing by $K$ the determinant formed with $a, h, g, \ldots, h, b, f, \ldots, g, f, c \ldots$ and by $A, H, G, \ldots, H, B, F, \ldots, G, F, C \ldots$ the inverse system of coefficients, we have
\[
K\frac{dX}{dx'} + A\delta a + H\delta h + G\delta g + \cdots + H\left(\frac{dM}{dx} - \frac{dL}{dy}\right) + G\left(\frac{dN}{dx} - \frac{dL}{dz}\right) + \cdots = 0,
\]
and similarly
\[
K\frac{dY}{dy'} + H\delta a + B\delta b + F\delta f + \cdots + H\left(\frac{dL}{dy} - \frac{dM}{dx}\right) + F\left(\frac{dN}{dy} - \frac{dM}{dz}\right) + \cdots = 0,
\]
\[
K\frac{dZ}{dz'} + G\delta g + F\delta f + C\delta c + \cdots + G\left(\frac{dL}{dz} - \frac{dN}{dx}\right) + F\left(\frac{dM}{dz} - \frac{dN}{dy}\right) + \cdots = 0.
\]

Hence, adding,
\[
K\left(\frac{dX}{dx'} + \frac{dY}{dy'} + \frac{dZ}{dz'} + \cdots\right) + A\delta a + B\delta b + C\delta c + \cdots + 2F\delta f + 2G\delta g + 2H\delta h + \cdots = 0;
\]
and thus we have as before, though with symbols bearing an entirely different signification,
\[
K\left(\frac{dX}{dx'} + \frac{dY}{dy'} + \frac{dZ}{dz'} + \cdots\right) + \delta K = 0,
\]
and thence $K\delta M - M\delta K = 0$, and $M = K$.

(The value of $K$ in this section may, I think, be conveniently termed ``the determinant of the \emph{vis viva},'' with respect to the variables $x, y, z, \ldots$. It may be remarked that ``the determinant of the \emph{vis viva}'' with respect to any other system of variables $u, v, w, \ldots$ is $= V^{-2} K$, $V$ as before.)

\subsection*{\S\,6. Third form of the equations of motion. [Hamiltonian form.]}

Here writing
\[
\frac{dT}{dx'} = \xi,\quad \frac{dT}{dy'} = \eta,\quad \ldots
\]
and taking $t, x, y, \ldots, \xi, \eta, \ldots$ for the variables of the problem [and considering $T$ to be expressed as a function of these variables: to denote this change it would have been proper to use instead of $T$ a new letter $H$,] the equations of motion reduce themselves to
\[
\frac{dx}{dt} = \frac{dT}{d\xi},\quad \frac{d\xi}{dt} = -\frac{dT}{dx} + P,
\]
\[
\frac{dy}{dt} = \frac{dT}{d\eta},\quad \frac{d\eta}{dt} = -\frac{dT}{dy} + Q,
\]
\[
\frac{dz}{dt} = \frac{dT}{d\zeta},\quad \frac{d\zeta}{dt} = -\frac{dT}{dz} + R,
\]
or putting for shortness
\[
\frac{dT}{d\xi} = \xi',\quad \frac{dT}{d\eta} = \eta',\quad \frac{dT}{d\zeta} = \zeta',
\]
\[
-\frac{dT}{dx} + P = X,\quad -\frac{dT}{dy} + Q = Y,\quad -\frac{dT}{dz} + R = Z,
\]
they become
\[
dt : dx : dy : dz : d\xi : d\eta : d\zeta \ldots
\]
\[
= 1 : \xi' : \eta' : \zeta' : X : Y : Z \ldots ,
\]
Hence writing the equation in $M$ under the form
\[
\delta M + M\left(\frac{d\xi'}{dx} + \frac{d\eta'}{dy} + \frac{d\zeta'}{dz} + \cdots + \frac{dX}{d\xi} + \frac{dY}{d\eta} + \frac{dZ}{d\zeta}+\cdots\right) = 0,
\]
\[
\Bigl(\text{where } \delta = \frac{d}{dt} + \xi'\frac{d}{dx} + \eta'\frac{d}{dy} + \cdots + X\frac{d}{d\xi} + Y\frac{d}{d\eta} + \cdots\Bigr),
\]
we see immediately that ($P, Q, \ldots$ being as before independent of the velocities, and consequently of $\xi, \eta, \zeta, \ldots$),
\[
\frac{d\xi'}{dx} + \frac{dX}{d\xi} = 0,\quad \frac{d\eta'}{dy} + \frac{dY}{d\eta} = 0,\quad \cdots
\]
Hence $\delta M = 0$, which is satisfied by $M = 1$.

\newpage

% ===========================================================================
% PDF page 307 -- Book page 285
% Paper [44]
% ===========================================================================
\section*{[44] On a Multiple Integral connected with the Theory of Attractions.}
\addcontentsline{toc}{section}{44. On a Multiple Integral connected with the Theory of Attractions}

\textit{[From the Cambridge and Dublin Mathematical Journal, vol.~ii.~(1847), pp.~219--223.]}

\bigskip

Mr Boole [in the Memoir ``On a Certain Multiple definite Integral'' Irish Acad.\ Trans.\ vol.~xxi.~(1848), pp.~40--150] has given for the integral with $n$ variables
\[
V = \int\!\!\int\cdots \frac{\phi(\alpha x + \beta y + \cdots)\,dx\,dy\cdots}{\bigl[\bigl(\tfrac{x-a}{f}\bigr)^2 + \bigl(\tfrac{y-b}{g}\bigr)^2 + \cdots\bigr]^{n/2 + q'}} \quad\cdots(1),
\]
limits
\[
\frac{x^2}{f^2} + \frac{y^2}{g^2} + \cdots = 1,
\]
the following formula, or one which may readily be reduced to that form,
\[
V = \frac{fg\cdots \pi^{n/2}}{\Gamma(\tfrac12 n + q')}\int_0^\infty \frac{S\,s^{-1-q'}\,ds}{\sqrt{R}\sqrt{s+f^2}\sqrt{s+g^2}\cdots} \quad\cdots(2),
\]
where
\[
S = \frac{(1-\sigma)^{-q'}}{\Gamma(1-q')}\int_0^1 t^{-q'}\phi\bigl[\sigma + t(1-\sigma)\bigr]\,dt \quad\cdots(3);
\]
\[
\sigma = \frac{a^2}{f^2 + s} + \frac{b^2}{g^2 + s} + \cdots - \frac{u^2}{\eta} \quad\cdots(4)
\]
and $\eta$ is determined by
\[
1 = \frac{a^2}{f^2 + \eta} + \frac{b^2}{g^2 + \eta} + \cdots - \frac{u^2}{\eta}.\footnotemark
\]
\footnotetext{See note at the end of this paper.}

Suppose $f = g = \cdots = \infty$; also assume
\[
\phi(z) = (\rho^2 + z^2)^{-\tfrac12 n - q} \quad\cdots(5),
\]
then the integral becomes
\[
U = \int\!\!\int\cdots \frac{dx\,dy\cdots}{(\alpha x^2 + \beta y^2 + \cdots + u^2)^{\tfrac12 n - q'}\,\{(x-a)^2 + (y-b)^2 + \cdots + u^2\}^{\tfrac12 n + q'}} \quad\cdots(6),
\]
the limits for each variable being $-\infty,\infty$.

Now, writing $f^2 s$ for $s$ and $f^2 \eta$ for $\eta$, the new value of $\eta$ reduces itself to zero, and
\[
U = \frac{\pi^{n/2}}{\Gamma(\tfrac12 n + q')}\int_0^\infty \frac{S\,ds}{s^{1+q'}}.
\]
Also $\sigma = 0$; but
\[
\phi z = \frac{1}{(\rho^2 + z^2)^{n/2 + q}},
\]
where $r^2 = a^2 + b^2 + \cdots$, whence also putting $-\frac{x^2}{s}$ for $t$, $\phi\{\sigma + t(1-\sigma)\}$ becomes
\[
\frac{1}{\{r^2 + s + t(1-\sigma) + u^2\}^{n/2 + q}},
\]
i.e.
\[
\frac{1}{(A + t)^{n/2 + q}}
\]
if for a moment
\[
A = \frac{r^2 + s + u^2}{1-\sigma}.
\]

Hence
\[
S = \frac{\Gamma(-q)}{\Gamma(\tfrac12 n + q + q')}\int_0^1 (t+A)^{-\tfrac12 n - q}\,s^{-q-1}\,ds,
\]
\[
\quad = \frac{\Gamma(\tfrac12 n + q + q')}{\Gamma(\tfrac12 n + q)}\,p,
\]
or substituting in $U$, and replacing $A$ by its value,
\[
U = \frac{\pi^{n/2}\,\Gamma(\tfrac12 n + q + q')}{\Gamma(\tfrac12 n + q)\,\Gamma(\tfrac12 n + q')}\int_0^\infty\!\int_0^\infty \frac{s^{-1-q'} t^{-1-q} dt\,ds}{\bigl(\sqrt{r^2 + s + u^2}\sqrt{s}\bigr)^{n + 2q + 2q'}},
\]
or, what comes to the same,
\[
U = \frac{\pi^{n/2}\,\Gamma(\tfrac12 n + q + q')}{\Gamma(\tfrac12 n + q)\,\Gamma(\tfrac12 n + q')}\int_0^\infty\int_0^\infty\frac{s^{-1-q'} t^{-1-q}\,ds\,dt}{J^{n/2 + q + q'}}
\]
where
\[
J = u^2 + r^2 + r^2 t.
\]

The only practicable case is that of $q' = -q$, for which
\[
U = \frac{\pi^{n/2}}{\Gamma(\tfrac12 n + q)\,\Gamma(\tfrac12 n - q)}\int_0^\infty\frac{B\,s^{-q-1}\,ds}{(u\sqrt{s} + \sqrt{r^2 + s + u^2})^n}.
\]
Consider the more general expression
\[
\theta = \int_0^\infty s^{-q-1}\phi\Bigl(\frac{u\sqrt{s} + \sqrt{r^2 + s + u^2}}{\sqrt{s}}\Bigr)\,ds \quad\cdots(9);
\]
by writing
\[
2u\sqrt{s} = \sqrt{(s' + 4uv)} \pm \sqrt{s'},
\]
the upper sign from $s = \infty$ to $s = \tfrac{r^2}{4u}$, and the lower one from $s = \tfrac{r^2}{4u}$ to $s = 0$, it is easy to derive
\[
\theta = (2u)^{2q}\int_0^\infty \frac{(\sqrt{s'} + 4uv) + \sqrt{s'}}{2\sqrt{s' + 4uv}}\,s'^{-q-1}\phi(s' + j + 2uv)\,ds' \quad\cdots(10).
\]

Now, by a formula which will presently be demonstrated,
\[
\frac{(\sqrt{s} + 4uv) + \sqrt{s}}{2\sqrt{s + 4uv}}\,s^{-q-1} = \frac{1}{\Gamma(-q)}\,e^{-s} \int_0^\infty s^{-1-q}(s + 4uv)^{-1-q}e^{-s}\,ds \quad\cdots(11);
\]
whence
\[
\frac{(\sqrt{s} + 4uv) + \sqrt{s}}{2\sqrt{s + 4uv}}\,\phi s = \frac{1}{\Gamma(-q)} \int_0^\infty s^{-1-q}(s + 4uv)^{-1-q}e^{-s}\phi(z + s)\,ds \quad\cdots(12);
\]
and hence in the particular case in question
\[
U = (2u)^{2q}\,\frac{\sqrt{\pi}\,\Gamma\bigl(\tfrac{n-1}{2}\bigr)}{2} \int_0^\infty s^{-1-q}(s + 4uv)^{-1-q}(s + j + 2uv)^{-1/2 + q}\,ds
\]
by means of the formula
\[
\Bigl(-\frac{d}{ds}\Bigr)^{1/2 - q}(s + 4uv)^{-q-1/2} = \frac{\Gamma(\tfrac12 - q)}{\Gamma(-q)}(s + 4uv)^{-1}.
\]
Thus, by merely changing the function,
\[
\theta = \frac{(2u)^{2q}}{\Gamma(-q)}\int_0^\infty s^{-1-q}(s + 4uv)^{-1-q}\Bigl(-\frac{d}{ds}\Bigr)^{-q}\phi(s + j + 2uv)\,ds.
\]

But as there may be some doubt about this formula, which is not exactly equivalent either to Liouville's or Peacock's expression for the general differential coefficient of a power, it is worth while to remark that, by first transforming the $\tfrac12 n$-th power into an exponential, and then reducing as above (thus avoiding the general differentiation), we should have obtained
\[
U = \frac{1}{\Gamma(-q)\,\Gamma(\tfrac12 n + q)\,\Gamma(\tfrac12 - q)}\int_0^\infty\!\int_0^\infty\!\int_0^\infty\cdots
\]
which reduces itself to the equation (14) by simply performing the integration with respect to $\theta$; thus establishing the formula beyond doubt\footnotemark. The integral may evidently be effected in finite terms when either $q$ or $q - \tfrac12$ is integral. Thus for instance in the simplest case of all, or when $q = -\tfrac12$,
\[
\frac{\sqrt{\pi}\,\Gamma\bigl(\tfrac{n+1}{2}\bigr)}{(j + 2uv)^{1/2 + n - 1}} = \int\!\!\int\cdots \frac{dx\,dy\cdots}{(x^2 + y^2 + \cdots + u^2)^{(n-1)/2}\,\{(x-a)^2 + \cdots + u^2\}^{(n+1)/2}},
\]
a formula of which several demonstrations have already been given in the Journal.
\footnotetext{A paper by M.\ Schl\"omilch ``Note sur la variation des constantes arbitraires d'une Int\'egrale d\'efinie,'' Crelle, t.~xxxiii.~[1846], pp.~268--280, will be found to contain formulae analogous to some of the preceding ones.}

The following is a demonstration, though an indirect one, of the formula (11): in the first place
\[
\int_{-\infty}^{\infty}\Bigl(\sqrt{s + 4uv} + \sqrt{s}\Bigr)\bigl(\sqrt{s + 4uv} - \sqrt{s}\bigr)^{-q-1}\,e^{ixx}\,dx
\]
\[
= \frac{2^{-q}}{\Gamma(\tfrac12 - q)\,uvi}\int_0^\infty (4u^2v^2 + x^2)^{q - 1/2}\,e^{ixx}\,dx \quad\cdots(16).
\]
(where as usual $i = \sqrt{-1}$): to prove this, we have
\[
\int_{-\infty}^{\infty}(4u^2v^2 + x^2)^{q-1/2}e^{ixx}\,dx = \frac{1}{\Gamma(\tfrac12 - q)}\int dx \int dt\, t^{-q-1/2}\,e^{-t(4u^2v^2 + x^2) + ixx}
\]
\[
= \frac{\sqrt{\pi}}{\Gamma(\tfrac12 - q)}\int_0^\infty dt\, t^{-q-1}\,e^{-4u^2v^2 t - x^2/(4t)}
\]
or, putting $iuv\sqrt{t} = \sqrt{(s + 4uv)} + \sqrt{s}$ (which is a transformation already employed in the present paper), the formula required follows immediately.

Now, by a formula due to M.\ Catalan, but first rigorously demonstrated by M.\ Serret (\emph{Liouville}, t.~viii.~[1843] p.~1), and by a slight modification in the form of this equation
\[
\int_{-\infty}^{\infty}(4u^2 v^2 + x^2)^{q-1/2}\,e^{ixx}\,dx = \frac{2^{-q}}{\Gamma(\tfrac12 - q)\,uvi}\int_0^\infty s^{-q-1/2}(s + 4uv)^{-q-1/2}\,e^{-s}\,ds,
\]
which, compared with (16), gives the required equation.

\bigskip
\textbf{Note.}---One of the intermediate formulae of Mr Boole [in the Memoir referred to] may be written as follows:
\[
S = \frac{1}{\pi}\int d\sigma \int d\nu\,\nu^{-q}\cos\bigl[(a - \sigma)z + \tfrac12 q\pi\bigr]\,\phi a,
\]
or what comes to the same thing, putting $t = \sqrt{-1}$, and rejecting the impossible part of the integral,
\[
S = -\frac{1}{2}e^{(q+1)\pi i/2}\int da \int dv\,v^{-q}\,e^{i(a - \sigma)v}\,\phi a,
\]
\[
S = \int I\,da\,za\,da,\quad I = -\frac12 e^{(q+1)\pi i/2}\int_0^\infty dv\,v^{-q}\,e^{i(a - \sigma)v}
\]
Now ($a - \sigma$) being positive,
\[
I = -\frac12 e^{(q+1)\pi i/2}\,\Gamma(q+1)\,e^{i(q+1)\pi/2}(a - \sigma)^{-q-1};
\]
or, retaining the real part only,
\[
I = \frac{\sin q\pi}{2}\,\Gamma(q+1)(a - \sigma)^{-q-1};
\]
i.e.
\[
I = -\frac{\pi}{2\,\Gamma(-q)}(a - \sigma)^{-q-1}.
\]
But ($a - \sigma$) being negative,
\[
I = -\tfrac12 e^{(q+1)\pi i/2}\,\Gamma(q+1)\,e^{-i(q+1)\pi/2}(\sigma - a)^{-q-1};
\]
\[
I = -\tfrac12\,e^{-q\pi i/2}\,\Gamma(q+1)(\sigma - a)^{-q-1},
\]
or, retaining the real part only, $I = 0$.

Hence
\[
S = -\frac{1}{2\Gamma(-q)}\int_{\sigma}^{a_1}(a - \sigma)^{-q-1}\,\phi a\,da,
\]
or putting $a = \sigma + t(1 - \sigma)$, or $a - \sigma = t(1 - \sigma)$,
\[
S = \frac{(1-\sigma)^{-q}}{\Gamma(1-q)}\int_0^1 t^{-q}\phi\bigl[\sigma + t(1 - \sigma)\bigr]\,dt:
\]
the expression in the text. Mr Boole's final value is
\[
S = \Bigl(-\frac{d}{d\sigma}\Bigr)^q\phi\sigma,
\]
which, though simpler, appears to me to be in some respects less convenient.

\newpage

% ===========================================================================
% PDF page 312 -- Book page 290
% Paper [45]
% ===========================================================================
\section*{[45] On the Theory of Elliptic Functions.}
\addcontentsline{toc}{section}{45. On the Theory of Elliptic Functions}

\textit{[From the Cambridge and Dublin Mathematical Journal, vol.~ii.~(1847), pp.~256--266.]}

\bigskip

Adopting the notation of the \emph{Fund.~Nova}, except that for shortness $\operatorname{sn} u(,1)$, $\operatorname{cn} u$, $\operatorname{dn} u$ are written instead of $\operatorname{sinam} u, \operatorname{cosam} u, \sqrt{\Delta\operatorname{am}\,u}$, let the functions $\Theta(u), H(u)$ be defined by the equations
\[
\Theta u = \sqrt{\Bigl(\frac{2Kk'}{\pi}\Bigr)}\,e^{-\int_0^u du \int_0^u k^2\operatorname{sn}^2 u\,du} \quad\cdots(1),
\]
\[
H u = -i e^{(\pi i u^2)/(4KK')}\Theta(u + iK') \quad\cdots(2),
\]
it is required from these equations to express $\operatorname{sn} u$ in terms of the functions $H(u)$, $\Theta(u)$. To accomplish this we have
\[
\frac{1}{\operatorname{sn} u}\frac{d}{du}\log\operatorname{sn} u = \operatorname{sn} u\,\frac{d}{du}\,\frac{1}{\operatorname{sn} u} - \frac{1}{\operatorname{sn}^2 u}\Bigl(\frac{d\operatorname{sn} u}{du}\Bigr)^2 \cdots
\]
\[
= -(1 + k^2) + 2k^2\operatorname{sn}^2 u - \frac{1}{\operatorname{sn}^2 u}\bigl[-(1 + k^2) + k^2\operatorname{sn}^4 u\bigr]
\]
\[
= k^2\operatorname{sn}^2 u - \frac{1}{\operatorname{sn}^2 u},
\]
whence also
\[
\frac{d^2}{du^2}\log\operatorname{sn} u = k^2\operatorname{sn}^2 u - k^2\operatorname{sn}^2(u + iK').
\]

If for a moment
\[
\psi_1 u = \int du\,\operatorname{sn}^2 u,\quad \psi_{11} u = \int du\,\psi_1 u,
\]
then
\[
\log\operatorname{sn} u = k^2\psi_{11} u - k^2\psi_{11}(u + iK') + Au + B;
\]
or writing $-u$ for $u$ and subtracting, $\psi_{11} u$ being an even function,
\[
2Au = i\pi - k^2\psi_{11}(iK' - u) + k^2\psi_{11}(iK' + u),
\]
or putting $u = K$,
\[
2AK = i\pi - k^2\psi_{11}(iK' - K) + k^2\psi_{11}(iK' + K).
\]
Now
\[
\operatorname{sn}^2(u + K) - \operatorname{sn}^2(u - K) = 0,
\]
and therefore $\psi_1(u + K) - \psi_1(u - K) = 2\psi_1 K$,
or
\[
\psi_{11}(iK' + K) - \psi_{11}(iK' - K) = 2iK'\psi_1 K.
\]
Also $E = K - k^2\psi_1 K$, i.e.\ $\psi_1 K = \frac{1}{k^2}\bigl(1 - \frac{E}{K}\bigr)$.

Hence
\[
A = iK'\Bigl(1 - \frac{E}{K}\Bigr) + \frac{\pi i}{2K},
\]
\[
\log\operatorname{sn} u = k^2\psi_{11} u - k^2\psi_{11}(u + iK') + uiK'\Bigl(1 - \frac{E}{K}\Bigr) + \frac{\pi i u}{2K} + B',
\]
i.e.
\[
\quad = k^2\psi_{11} u - k^2\psi_{11}(u + iK') + \tfrac12\bigl[(u + iK')^2 - u^2\bigr]\Bigl(\frac{iK'}{KK'} - \frac{1}{k^2}\Bigr) + \frac{\pi i u}{2K} + B',
\]
\[
\log\operatorname{sn} u = \log\Theta(u + iK') - \log\Theta u + \frac{\pi i u}{2K} + B',
\]
or, changing the constant,
\[
\operatorname{sn} u = C e^{\pi i u/(2K)}\,\frac{\Theta(u + iK')}{\Theta u}.
\]

Now, to determine $C$, write $u - iK'$ for $u$; this gives
\[
\frac{1}{k\operatorname{sn} u} = \frac{C e^{\pi i(u - iK')/(2K)}\,\Theta(u - iK')}{\Theta(u - iK')},
\]
and again changing $u$ into $-u$,
\[
-\operatorname{sn} u = C e^{-\pi i u/(2K)}\,\frac{\Theta(-u + iK')}{\Theta(-u)};
\]
whence, multiplying these last two equations,
\[
C^2 = \frac{1}{k}\,e^{-\pi K'/(2K)};
\]
or
\[
C = \frac{1}{\sqrt{k}}\,e^{-\pi K'/(4K)},
\]
whence
\[
\operatorname{sn} u = \frac{1}{\sqrt{k}}\,e^{\pi i u/(2K) - \pi K'/(4K)}\,\frac{\Theta(u + iK')}{\Theta u}
\]
i.e.
\[
\sqrt{k}\operatorname{sn} u = \frac{H(u)}{\Theta(u)} \quad\cdots(3);
\]
and the equations (1), (2) and (3) may be considered as comprehending the theory of the functions $H(u), \Theta(u)$. The preceding process is, in fact, the converse of that made use of in the \emph{Fund.~Nova}; Jacobi having obtained for $\operatorname{sn} u$ an expression in the form of a fraction, takes the numerator of it for $H(u)$ and the denominator for $\Theta(u)$, and thence deduces the equations (1), (2), the intermediate steps of the demonstration being conducted by means of infinite series; the necessity of which is avoided by the preceding investigation.

I proceed to investigate certain results relating to these functions, and to the theory of elliptic functions which have been given by Jacobi in two papers, ``Suite des notices sur les fonctions elliptiques,'' Crelle, t.~iii.~[1828] p.~306, and t.~iv.~(4)~[1829] p.~185, but without demonstration.

In the first place, the equation
\[
\frac{d^2 \Sigma}{du^2} - 2 k^2\Bigl(\frac{E}{K} - \frac{1}{k^2}\Bigr)\Sigma + 2k^2\frac{dk}{du}\Sigma = 0 \quad\cdots(4),
\]
is satisfied by $\Sigma = \Theta(u)$ or $\Sigma = H(u)$. It will be sufficient to prove this for $\Sigma = \Theta(u)$, since a similar demonstration may easily be found for the other value. The following preliminary formulae will be required:
\[
k\frac{dE}{dk} = E - K,\quad \frac{dK}{dk} = \frac{E - k'^2 K}{k k'^2},\quad K' = \frac{dK}{dk'} = -\frac{k}{k'},\quad KK' - EK' - E'K = -\tfrac12\pi,
\]
which are all of them known.

Now, writing $\Theta(u)$ under the slightly more convenient form
\[
\Theta u = \sqrt{\frac{2Kk'}{\pi}}\,e^{-\int_0^u du \int_0^u dn^2 u}
\]
we have
\[
\frac{d\Theta u}{du} = \Bigl(-\int_0^u du\,\operatorname{dn}^2 u - \frac{2K}{k k'}u\Bigr)\Theta u = u\Bigl(k'^2 - \frac{E}{K}\Bigr)\Theta u + \int_0^u du\,k^2\operatorname{cn}^2 u\,\Theta u,
\]
\[
\frac{d^2\Theta u}{du^2} = \Bigl(k'^2 - \frac{E}{K}\Bigr)\Theta u + \cdots,
\]
\[
\frac{d^2\Theta u}{du^2} = \operatorname{dn}^2 u - \frac{E}{K} + \tfrac12 u\,k^2\,\Theta u\cdots
\]

The success of the process depends upon a transformation of the double integral
\[
\int_0^u du \int_0^u du \frac{d}{dk}\operatorname{dn}^2 u:
\]
to effect this we have
\[
\frac{d}{dk}\operatorname{dn}^2 u = -2k\operatorname{sn} u\Bigl(\operatorname{sn} u + k\frac{d}{dk}\operatorname{sn} u\Bigr);
\]
but, by a known formula,
\[
k^2\frac{d}{dk}\operatorname{sn} u = -k\operatorname{cn} u\,\operatorname{dn} u \int_0^u\operatorname{cn}^2 u\,du + k\operatorname{cn}^2 u\,\operatorname{sn} u;
\]
Hence
\[
\operatorname{sn} u + k\frac{d}{dk}\operatorname{sn} u = \frac{1}{k k'^2}\operatorname{sn} u\,\operatorname{dn}^2 u - \frac{1}{k^2}\operatorname{cn} u\,\operatorname{dn} u \int_0^u du\,\operatorname{cn}^2 u,
\]
\[
\frac{d}{dk}\operatorname{dn}^2 u = -\frac{2k}{k k'^2}\bigl(\operatorname{sn}^2 u\,\operatorname{dn}^2 u - k^2\operatorname{sn} u\,\operatorname{cn} u\,\operatorname{dn} u \int_0^u du\,\operatorname{cn}^2 u\bigr)
\]
\[
\quad = -\frac{1}{k k'^2}\operatorname{sn}^2 u\,\operatorname{dn}^2 u + \tfrac12 k^2\Bigl[\frac{d}{du}\operatorname{cn}^2 u\int du\,\operatorname{cn}^2 u\Bigr];
\]
whence
\[
\int_0^u du \int_0^u du\,\frac{d}{dk}\operatorname{dn}^2 u = -\tfrac{1}{2}\Bigl\{\int_0^u du \int_0^u du\,\operatorname{sn}^2 u\,\operatorname{dn}^2 u + \tfrac12 k^2\int_0^u du\,\bigl(\operatorname{cn}^2 u \int du\,\operatorname{cn}^2 u\bigr)\Bigr\}
\]
\[
= -\frac{1}{k k'^2}\bigl\{\int_0^u du\,\int_0^u du\,(2\operatorname{sn}^2 u\,\operatorname{dn}^2 u - k^2\operatorname{cn}^4 u) + \tfrac12 k^2\bigl(\int_0^u du\,\operatorname{cn}^2 u\bigr)^2\bigr\}.
\]
But
\[
\frac{d^2}{du^2}\operatorname{sn}^2 u = 2(\operatorname{cn}^2 u\,\operatorname{dn}^2 u - \operatorname{sn}^2 u\,\operatorname{dn}^2 u - k^2\operatorname{sn}^2 u\,\operatorname{cn}^2 u) = 2(k'^2 - 2\operatorname{sn}^2 u\,\operatorname{dn}^2 u + k^2\operatorname{cn}^4 u);
\]
or, integrating,
\[
\operatorname{sn}^2 u = k'^2 u^2 - 2\int_0^u du \int_0^u du\,(2\operatorname{sn}^2 u\,\operatorname{dn}^2 u - k^2\operatorname{cn}^4 u);
\]
whence at length
\[
\int_0^u du \int_0^u du\,\frac{d}{dk}\operatorname{dn}^2 u = -\tfrac12 k u^2 + \frac{k}{k k'^2}\operatorname{sn}^2 u + \tfrac12 k^2\Bigl(\int_0^u du\,\operatorname{cn}^2 u\Bigr)^2.
\]

Also
\[
\frac{d}{dk}\frac{E - K}{k k'^2} = \frac{1}{k k'^2}\frac{dE}{dk} - \frac{1}{k k'^2}\frac{dK}{dk} - \cdots
\]
so that
\[
\frac{d}{dk}\frac{E - K}{k k'^2} = \tfrac12 k^2\bigl[k - \tfrac12 k^3 - k^3 - 3 - k k^3\cdots\bigr].
\]

Substituting these values of $\frac{d\Theta u}{du},\ \frac{d^2\Theta u}{du^2},\ \frac{d\Theta u}{dk},\ \Theta u$ in the equation (4) in the place of the corresponding differential coefficients of $\Sigma$, all the terms vanish, or the equation is satisfied by $\Sigma = \Theta(u)$, and similarly it would be satisfied by $\Sigma = H(u)$.

Assume now
\[
u = \frac{2K}{\pi}\,\omega,\quad \frac{K'}{K} = \frac{1}{\pi}\,\Omega.
\]
then observing the equation
\[
\frac{d}{du} = \frac{\pi}{2K}\frac{d}{d\omega},\quad \frac{d^2}{du^2} = \frac{\pi^2}{4K^2}\frac{d^2}{d\omega^2},
\]
\[
\frac{d}{du} = u\Bigl(k'^2 - \frac{E}{K}\Bigr)\frac{d}{dv} - \frac{2K}{\pi k k'^2}\frac{d}{d\Omega},
\]
whence, substituting in the equation (4), this becomes
\[
\frac{d^2\Sigma}{dv^2} - 4\,\frac{d\Sigma}{d\Omega} = 0 \quad\cdots(5)
\]
which is of course satisfied as before by $\Sigma = \Theta(u)$, or $\Sigma = H(u)$, an equation demonstrated in a different manner (by means of expansions) by Jacobi in the Memoirs referred to.

Consider next the equation
\[
\frac{d^2\Sigma}{du^2} - 2n k^2\Bigl(\frac{E}{K} - \frac{1}{k^2}\Bigr)\Sigma + 2nk^2\frac{dk}{du}\Sigma = 0 \quad\cdots(6),
\]
($n$ being any positive integer number). Then, by assuming
\[
\omega = n\,\frac{\pi K'}{2K},\quad v = \frac{n\pi u}{2K},
\]
we should be led as before to the equation (6) modulus satisfied (corresponding to (5). Hence assuming a transformation of the $n$-th order; then $\lambda, \lambda'$ being the complete functions corresponding to this modulus,
\[
\lambda' = \frac{n K'}{K},\text{ so that the equation (6) will be satisfied by assuming }\Sigma = \Theta_1(nu) \text{ or } \Sigma = H_1(nu),
\]
where $\Theta_1, H_1$ correspond to the new modulus $\lambda$.

Assume now in the equation (6),
\[
\Sigma = \Bigl(\frac{2}{\pi}\Bigr)^{1/2}\Bigl(\frac{n-1}{2}\Bigr)\cdots
\]
Hence, substituting,
\[
-\Theta^n u \cdot z - 2nu\bigl(k'^2 - \tfrac{E}{K}\bigr)\frac{d}{du}\bigl(\Theta^n u \cdot z\bigr) + 2nk^2 k'\bigl(K k'\bigr)^{-\tfrac12 (n-1)}\frac{d}{dk}\bigl[(K k')^{-\tfrac12(n-1)}\Theta^n u \cdot z\bigr] = 0;
\]
but
\[
\frac{d}{du}(\Theta^n u\cdot z) = n\Theta^{n-1} u\,\frac{d\Theta u}{du}\,z + \Theta^n u\,\frac{dz}{du},
\]
or effecting the differentiation, and eliminating $\frac{d\Theta u}{du}$ by means of the equation obtained from (4) by writing $\Sigma = \Theta u$,
\[
(K k')^{-\tfrac12(n-1)}\frac{d}{dk}\bigl[(K k')^{-\tfrac12(n-1)}\Theta^n u \cdot z\bigr]
\]
\[
= \Theta^n u\,\frac{dz}{dk} - nz\,\frac{d^2\Theta u}{du^2} - 2\frac{d}{du}\Bigl(1 - \frac{E}{K}\Bigr)\,\frac{d\Theta u}{du} - \frac{n-1}{2K k k'}\,\Theta^n u\,z.
\]

Substituting in (6) and reducing,
\[
\frac{d^2z}{du^2} + n(n-1)\,k^2\,\operatorname{sn}^2 u\cdot z = 0 \quad\cdots(7);
\]
\[
\frac{d^2\log\Theta u}{du^2} = \frac{1}{k k'^2}\Bigl(k - \frac{E}{K}\Bigr) = k^2\int du\,\operatorname{dn}^2 u,
\]
\[
\frac{d^2\log\Theta u}{du^2} + 2nk^2 \int du\,\operatorname{cn}^2 u\,\frac{d z}{du} + 2n k k'\,\frac{dz}{dk} + n(n-1)\,k^2\,\operatorname{sn}^2 u \cdot z = 0,
\]
which is therefore satisfied by
\[
z = \frac{\sqrt{\lambda}\operatorname{sn}_1(nu)}{(2 \cdot K k')^{(n-1)/2}}\cdot\frac{1}{\Theta^n u} = \Bigl(\frac{2}{\pi}\Bigr)^{(n-1)/2}\,\frac{H_1(nu)}{\Theta^n u};
\]
and each of these values is an algebraical function of $\operatorname{sn} u$, (viz.\ either a rational function or a rational function multiplied by $\operatorname{cn} u\operatorname{dn} u$). Also, in the transformation of the $n$-th order,
\[
\frac{H_1(nu)}{\Theta^n u} = \frac{\Theta_1(nu)}{\Theta^n u},
\]
so that it is clear that the above values of $z$ may be taken for the denominator and numerator respectively of $\sqrt{\lambda}\operatorname{sn}_1 nu$; i.e.\ these quantities each of them satisfy the equation (7).

By assuming
\[
z = \sqrt{\lambda}\operatorname{sn} u,\quad x = \sqrt{k}\,\operatorname{sn} u
\]
this becomes
\[
\frac{d^2z}{dx^2} - n(n-1)x^2 z + (n-1)(\alpha x - 2x^3)\frac{dz}{dx} + (1 - \alpha x^2 + x^4)\frac{d^2 z}{dx^2} - 2n(a^2 - 4)\,\frac{dz}{dx} = 0 \quad\cdots(8);
\]
which is therefore satisfied by assuming for $z$ either the numerator or the denominator of $\sqrt{\lambda}\operatorname{sn}_1 u$ (the transformation of the $n$-th order), which is the form in which the property is given by Jacobi.

In the case where $n$ is odd, the denominator is of the form
\[
B_{n-1} - 1 - B_3 x^4 \cdots + B_{n-1} x^{n-1},
\]
and then the numerator is $Bx$,
where
\[
B = x(B_{n-1} \cdots + B_2 x^{n-3} + B_0 x^{n-1}),
\]
\[
B_{n-1} = \sqrt{\frac{\lambda}{nk}},\quad B_{n-1} = \sqrt{\frac{n\lambda}{k^n}},
\]
and all the remaining coefficients may be determined from these, the modular equation being supposed known. But the principal use of the formula is for the multiplication of elliptic functions, which it is well known corresponds to the case where $n$ is a square number. Writing $n = \nu^2$, when $\nu$ is odd, the denominator is
\[
1 + B_2 x^4 \cdots + B_{\nu^2 - 3} x^{\nu^2 - 5} \pm \nu x^{\nu^2 - 1},
\]
(the $\pm$ sign according as $\nu = (4p + 1)$ or $(4p - 1)$); and the numerator is obtained from this by multiplying by $x$ and reversing the order of the coefficients. When $\nu$ is even the denominator is
\[
1 + B_2 x^4 \cdots + B_{\nu^2 / 2 - 1} \pm x^{\nu^2},
\]
($+$ or $-$, according as $\nu = 4p$ or $\nu = 4p + 2$), so that there are only half as many coefficients to be determined; but then the numerator must be separately investigated.

In general, by leaving $n$ indeterminate, and integrating in the form of a series arranged according to ascending powers of $x^2$; then, whenever $n$ is a square number, the series terminates and gives the denominator of the corresponding formula of multiplication; but the general form of the coefficients has not hitherto been discovered.

By writing $\frac{1}{x}$ instead of $x$, and then making $n$ infinite, the equation (8) takes the form
\[
\frac{d^2z}{dx^2} + \frac{1}{x}\frac{dz}{dx} - 2(\alpha x - 2x^3)\frac{dz}{dx} = 0 \quad\cdots(8'),
\]
and it is worth while, before attempting the solution of the general case, to discuss this more simple one\footnotemark.
\footnotetext{Writing $(\beta + 2)$ for $\alpha$, and putting $z = e^{\beta x^2}\,\rho$, this becomes
\[
\frac{d^2\rho}{dx^2} = -\frac{1}{x}\frac{d\rho}{dx} + \bigl[\beta(\beta + 2)x^2\bigr]\rho;
\]
and if $\rho = \Sigma Z_n x^{4n}$,
\[
(8n + 1)Z_n = \bigl(\tfrac14 \beta^2 + 2n - 2 - x^2\bigr)Z_{n-1};
\]
from which the successive values of $Z_0, Z_1$, \&c.\ might be calculated.}

Assume
\[
z = 1 + C_1\,\frac{x^2}{1\cdot 2}\,a + \cdots + C_r\,\frac{x^{2r}}{1\cdot 2 \cdots 2r}\,a^r \cdots,
\]
then it is easy to obtain
\[
C_{r+1} = -(2r+1)(2r+2)C_r - (2r+2)\alpha C_{r+1} + 2(\alpha^2 - 4)C_r'.
\]
The general form may be seen to be
\[
C_r = (-)^{r+1}\bigl[2^{r-1} {}_0 C_r^1 a^{r-1} + 2^{r-2} {}_0 C_r^2 a^{r-2} + \cdots\bigr],
\]
and then
\[
C_{r+1}^p - p C_r^p = -r(2r-1)C_{r-1}^p + 16(r + 2 - 2p) C_{r-1}^{p-1}.
\]

The complete value of $C_r^p$ (assuming $C_1^p = 0$) is given by an equation of the form
\[
C_r^p = {}^0 C_r^p + {}^1 C_r^p \cdot 2^r + {}^2 C_r^p \cdot 3^r \ldots + {}^{p-1} C_r^p \cdot p^r,
\]
where ${}^0C_r^p, {}^1 C_r^p$ are algebraical functions of $r$ of the degrees $2p-2, 2p-4$, \&c.\ respectively; but as I am not able completely to effect the integration, and my only object being to give an idea of the law of the successive terms, it will be sufficient to consider the first or algebraical term ${}^0C_r^p$, which is determined by the same equation as $C_r^p$, and is moreover completely determined by this equation and the single additional relation $C_1^p = 1$, since the arbitrary constants of the integration affect only the terms multiplied by $2^r, 3^r$, \&c.

Assume ${}^0C_r^p =$
\[
-\tfrac12\bigl\{2^{p-1}L^p\,[r-2]^{2p-2} + 2^{p-3}M^p\,[r-3]^{2p-1} + \cdots 2^{-p+1}X_p [r - 2p]^{2p}\bigr\};
\]
and substituting this value,
\[
(1-p)L^p = (1-p)(L^{p-1}),
\]
\[
(1-p)M^p - 2p(2-2p)L^p = (1-p)(M^{p-1} - 11^{p-1}),
\]
\[
(1-p)N^p - 2p(3-2p)M^p = (1-p)(N^{p-1} - 7M^{p-1} + 12L^{p-1}),
\]
\[
(1-p)O^p - 2p(4-2p)N^p = (1-p)(O^{p-1} - 3N^{p-1} + 30M^{p-1}),
\]
the law of which is obvious, the coefficients on the second side in the $q$-th line being $1, 4q - 19$, and $(2q - 3)(2q - 2)$ respectively. By successive integrations and substitutions
\[
L^p - L^{p-1} = 0,\quad L^p = 1,
\]
\[
M^p - M^{p-1} = 4p - 11,\quad M^p = (p - 1)(2p - 7),
\]
\[
N^p - N^{p-1} = -8p^3 + 26p^2 + 49p - 114;
\]
(the constants determined by $M^1 = 0, N^1 = 0, O^1 = 0, P^1 = 0,\ldots$ so as to make $C_r^p$ contain positive powers only of $r$).

The following are a few of the complete values of $C_r^p$, the constants determined so as to satisfy $C_{p+1}^p = 0$ (except $C_2^1 = 1$), and the factorials being partially developed in powers of $r$, viz.
\[
C_r^1 = 1,
\]
\[
C_r^2 = (r - 3)(2r - 7),
\]
\[
C_r^3 = \tfrac14(r - 4)(r - 5)(4r^2 - 24r + 51),
\]
\[
C_r^4 = \tfrac14\bigl\{(r - 5)(r - 6)(r - 7)(8r^3 - 60r^2 + 286r + 63) + 384(9r^2 - 93r + 242 - 2\cdot 4^{r-5})\bigr\},
\]
\&c.\
\[
\text{(it is curious that }C_2^4, C_3^4, C_4^4\text{ all three of them vanish).}
\]
It seems hopeless to continue this investigation any further.

Returning to the equation (8), and assuming for $z$ an expression of the same form as before for the coefficients $C_r$, we have, corresponding to the equations before found for the coefficients,
\[
C_{r+1} = -(2r + 1)(2r + 2)(n - 2r)(n - 2r - 1)C_r - (2r + 2)(n - 2r - 2)\alpha\,C_{r+1} + 2n(\alpha^2 - 4)\,C_r'.
\]
The case corresponding to the denominator in the multiplication of elliptic functions is that of $C_0 = 1, C_1 = 0$. It is easy to form the table---
\[
C_0 = 1,
\]
\[
C_1 = 0,
\]
\[
C_2 = -2n(n - 1),
\]
\[
C_3 = 8n(n - 1)(n - 4)\alpha,
\]
\[
C_4 = -4n(n - 1)(n - 4)[n + 75] - 32n(n - 1)(n - 4)(n - 9)\alpha^2,
\]
\[
C_5 = 96n(n - 1)(n - 4)(n - 9)[n + 44]\alpha + 128n(n - 1)(n - 4)(n - 9)(n - 16)\alpha^3,
\]
\[
C_6 = -24n(n - 1)(n - 4)(n - 9)[17n^2 + 403n + 9000]
\]
\[
\quad - 960n(n - 1)(n - 4)(n - 9)(n - 16)[n + 41]\alpha^2
\]
\[
\quad - 512n(n - 1)(n - 4)(n - 9)(n - 16)(n - 25)\alpha^4,
\]
\[
C_7 = +96n(n - 1)(n - 4)(n - 9)(n - 16)[79n^2 + 2825n + 36180]\alpha
\]
\[
\quad + 7168n(n - 1)(n - 4)(n - 9)(n - 16)(n - 25)[n + 42]\alpha^3
\]
\[
\quad + 2048n(n - 1)(n - 4)(n - 9)(n - 16)(n - 25)(n - 36)\alpha^5,
\]
\[
C_8 = -48n(n - 1)(n - 4)(n - 9)[283 n^4 - 26978 n^3 + 277827 n^2 - 5491932 n + 127764000]
\]
\[
\quad - 3840 n(n - 1)(n - 4)(n - 9)(n - 16)(n - 25)[23 n^2 + 1069 n + 23436]\alpha^2
\]
\[
\quad - 15360 n(n - 1)(n - 4)(n - 9)(n - 16)(n - 25)(n - 36)[3n + 133]\alpha^4
\]
\[
\quad - 8192 n(n - 1)(n - 4)(n - 9)(n - 16)(n - 25)(n - 36)(n - 49)\alpha^6,
\]
\&c.,
in which of course the coefficient of the highest power of $n$, in the successive coefficients $C_r$, is the value of $C_r$ obtained from the equation (8'). With regard to the law of these coefficients I have found that
\[
C_r = (-)^{r+1}\bigl\{2^{r-1} n(n - 1^2)\cdots\{n - (r - 1)^2\}\,C_r' \alpha^{r-1}
\]
\[
\qquad\qquad + 2^{r-2} n(n - 1^2)\cdots\{n - (r - 2)^2\}\,C_r'' \alpha^{r-2}
\]
\[
\qquad\qquad + 2^{r-3} \cdot 1 \cdot n(n - 1^2)\cdots\{n - (r - 3)^2\}\,C_r''' \alpha^{r-3}+\cdots\bigr\}
\]
(where however the next term does not contain, as would at first sight be supposed, the factor $n(n - 1^2)\cdots\{n - (r - 4)^2\}$). And then
\[
C_1' = 1,
\]
\[
C_2' = (r - 3)\bigl[n(2r - 7) + (r - 1)(8r - 7)\bigr],
\]
\[
C_3' = \tfrac14(r - 4)(r - 5)\bigl[n^2(4r^2 - 24r + 51)
\]
\[
\quad + n(32r^3 - 220r^2 + 412r - 255)
\]
\[
\quad + 2(r - 1)(r - 2)(32r^2 - 88r + 51)\bigr].
\]

In conclusion may be given the following results, in which, recapitulating the notation
\[
x = \tfrac{i}{2}\sqrt{k}\operatorname{sn} u,\quad a = k + \tfrac{1}{k},\quad \beta = \sqrt{(1 - \alpha x^2 + x^4)},
\]
\[
\sqrt{k}\operatorname{sn} 2u = \frac{2x\beta}{1 - x^4},
\]
\[
\sqrt{k}\operatorname{sn} 4u = \frac{4x\beta(1 - x^4)(1 - 6x^4 + x^8 - 4\alpha x^2 + 4\alpha x^6)}{1 - 20x^4 + (26 + 160\alpha^2)x^8 + 32\alpha x^{10} - 20x^{12} + x^{16}},
\]
\[
\sqrt{k}\operatorname{sn} 5u = x\bigl\{5 - 20\alpha x^2 + (62 + 16\alpha^2)x^4 - 80\alpha x^6 - 105 x^8 + 360\alpha x^{10} - (300 + 240\alpha^2)x^{12}
\]
\[
\qquad\qquad + (368\alpha + 64\alpha^3)x^{14} - (125 + 160\alpha^2)x^{16} + 140 \alpha x^{18} - 50 x^{20} + x^{24}\bigr\}
\]
\[
\div\bigl\{1 - 50 x^4 + 140\alpha x^6 - (125 + 160\alpha^2)x^8 + (368\alpha + 64\alpha^3)x^{10} - (300 + 240\alpha^2)x^{12}
\]
\[
\qquad\qquad + 360\alpha x^{14} - 105 x^{16} - 80 \alpha x^{18} + (62 + 16\alpha^2)x^{20} - 20\alpha x^{22} + 5 x^{24}\bigr\},
\]
\&c.

Thus, writing $-x^2$ for $x^2$, $k = 1$, and therefore $\alpha = 2$,
\[
\tan 3u = \frac{x(3 - 6x^2 + 6x^4 - x^6) - x(3 - x^2)(1 + x^2)^2}{8x^2 - 3x^6} = \frac{x(3 - x^2)}{1 - 3x^2},
\]
where $x = \tan u$. (And in general in reducing $\tan nu$ the extraneous factor in the numerator and denominator is $(1 + x^2)^{(n^2 - 1)/(n - 1)}$.)

\newpage

% ===========================================================================
% PDF page 323 -- Book page 301
% Paper [46]
% ===========================================================================
\section*{[46] Note on a System of Imaginaries.}
\addcontentsline{toc}{section}{46. Note on a System of Imaginaries}

\textit{[From the Philosophical Magazine, vol.~xxx.~(1847), pp.~257--258.]}

\bigskip

The octuple system of imaginary quantities $i_1, i_2, i_3, i_4, i_5, i_6, i_7$, which I mentioned in a former paper [21], (and the conditions for the combination of which are contained in the symbols
\[
123,\ 246,\ 374,\ 145,\ 275,\ 365,\ 167,
\]
i.e.\ in the formulae
\[
i_1 i_2 = i_3,\ i_3 i_2 = -i_1,\ i_1 i_3 = -i_2,\ i_2 i_1 = -i_3,\ i_3 i_1 = i_2,\ i_2 i_3 = i_1,
\]
with corresponding formulae for the other triplets $246$, \&c.,) possesses the following property; namely, if $i_\alpha, i_\beta, i_\gamma$ be any three of the seven quantities which do not form a triplet, then
\[
(i_\alpha i_\beta)\cdot i_\gamma = -i_\alpha\cdot(i_\beta i_\gamma).
\]
Thus, for instance,
\[
(i_2 i_4)\cdot i_5 = -i_6 \cdot i_5 = -i_2,
\]
but
\[
i_2\cdot(i_4 i_5) = i_2\cdot i_1 = i_3,
\]
and similarly for any other such combination. When $i_\alpha, i_\beta, i_\gamma$ form a triplet, the two products are equal, and reduce themselves each to $-1$, or each to $+1$, according to the order of the three quantities forming the triplet. Hence in the octuple system in question neither the commutative nor the distributive law holds, which is a still wider departure from the laws of ordinary algebra than that which is presented by Sir W.\ Hamilton's quaternions.

I may mention, that a system of coefficients, which I have obtained for the rectangular transformation of coordinates in $n$ dimensions (Crelle, t.~xxxii.~[1846] ``Sur quelques propri\'et\'es des D\'eterminans gauches'' [52]), does not appear to be at all connected with any system of imaginary quantities, though coinciding in the case of $n = 3$ with those mentioned in my paper ``On Certain Results relating to Quaternions,'' \emph{Phil.\ Mag.} Feb.\ 1845, [20].

\newpage

% ===========================================================================
% PDF page 324 -- Book page 302
% Paper [47]
% ===========================================================================
\selectlanguage{french}
\section*{[47] Sur la Surface des Ondes.}
\addcontentsline{toc}{section}{47. Sur la Surface des Ondes}

\textit{[From the Journal de Math\'ematiques Pures et Appliqu\'ees (Liouville), tom.~xi.~(1846), pp.~291--296.]}

\bigskip

La surface des ondes est un cas particulier d'une autre surface \`a laquelle on peut donner le nom de \emph{t\'etra\'edroide}, et dont voici la propri\'et\'e fondamentale:

``Le t\'etra\'edroide est une surface du quatri\`eme ordre, qui est coup\'ee par les plans d'un certain t\'etra\`edre suivant des paires de coniques par rapport auxquelles les trois sommets du t\'etra\`edre dans ce plan sont des points conjugu\'es. De plus: les seize points d'intersection des quatre paires de coniques sont des points singuliers de la surface, c'est-\`a-dire des points o\`u, au lieu d'un plan tangent, il y a un c\^one tangent du second ordre.''

On verra plus loin que la surface est d\'efinie suffisamment en disant qu'un seul de ces points \'etait un point singulier; l'existence des quinze autres points singuliers est donc une propri\'et\'e assez remarquable. Mais, avant d'aller plus loin, il convient de rappeler la signification des points conjugu\'es par rapport \`a une conique: cela veut dire que la polaire de chacun des trois points passe par les deux autres. Parmi les propri\'et\'es d'un tel syst\`eme, on peut citer celle-ci:

``Les points d'intersection de deux coniques, par rapport auxquelles les m\^emes trois points sont des points conjugu\'es, sont situ\'es deux \`a deux sur six droites, lesquelles passent deux \`a deux par ces trois points.''

De l\`a: ``Les quatre points singuliers compris dans chaque face du t\'etra\`edre sont situ\'es deux \`a deux sur trois paires de droites, lesquelles passent par les trois sommets dans cette face.''

Autrement dit, les seize points singuliers sont situ\'es deux \`a deux sur vingt-quatre droites, lesquelles passent six \`a six par les sommets et sont situ\'ees six \`a six dans les plans du t\'etra\`edre. Donc, en consid\'erant les six droites qui passent par le m\^eme sommet, par ces droites combin\'ees deux \`a deux, il passe quinze plans, y compris trois plans du t\'etra\`edre; en excluant ceux-ci, il y a pour chaque sommet douze plans, chacun desquels passe par quatre points singuliers; donc la courbe d'intersection d'un de ces plans avec la surface a quatre points doubles, ce qui ne peut pas arriver pour une courbe du quatri\`eme ordre, \`a moins qu'elle ne se r\'eduise \`a deux coniques.

Donc: ``Il y a, outre les plans du t\'etra\`edre, quarante-huit plans, douze par chaque sommet, lesquels rencontrent la surface suivant des paires de coniques.''

On pourrait de m\^eme chercher combien il y a de plans qui rencontrent la surface suivant des courbes avec trois points doubles, \&c. Passons aux c\^ones circonscrits; on d\'emontre facilement par l'analyse cette propri\'et\'e:

``Les c\^ones circonscrits \`a la surface ayant pour sommets les sommets du t\'etra\`edre, se r\'eduisent \`a des paires de c\^ones du second ordre, lesquelles la touchent suivant les deux coniques de la face oppos\'ee.''

De plus: ``Les seize plans qui touchent quatre \`a quatre ces paires de c\^ones sont plans tangents singuliers, dont chacun touche la surface suivant une conique.''

Il y a de plus, entre ces plans, des relations analogues \`a celles qui existent entre les points singuliers, de mani\`ere que l'on d\'eduit de m\^eme le th\'eor\`eme:

``Il y a quarante-huit points, douze \`a douze dans les quatre plans du t\'etra\`edre, pour chacun desquels les c\^ones circonscrits se r\'eduisent \`a des paires de c\^ones du second ordre.''

Ajoutons que ces quarante-huit points correspondent d'une mani\`ere particuli\`ere aux quarante-huit plans, et que les c\^ones circonscrits touchent la surface suivant les coniques situ\'ees dans les plans correspondants.

La r\'eciproque d'une surface du quatri\`eme ordre est, en g\'en\'eral, de l'ordre trente-six; mais ici, \`a cause des seize points singuliers, cet ordre se r\'eduit de trente-deux, savoir, \`a quatre. Et les propri\'et\'es qui viennent d'\^etre \'enonc\'ees par rapport aux c\^ones circonscrits montrent que la r\'eciproque \'etant de cet ordre est n\'ecessairement une surface de la m\^eme esp\`ece; c'est-\`a-dire:

``La r\'eciproque d'un t\'etra\'edroide est aussi un t\'etra\'edroide.''

Donc le t\'etra\'edroide est surface de la quatri\`eme classe. Puisque cette surface n'a pas de lignes doubles ou de lignes de rebroussement, il n'y a pas de r\'eduction dans le nombre qui exprime le rang de la surface, lequel est ainsi \'egal \`a douze, c'est-\`a-dire le c\^one circonscrit est ordinairement de l'ordre douze. Mais nous venons de voir que ce c\^one est seulement de la quatri\`eme classe (en effet, la classe du c\^one est la m\^eme chose que celle de la surface); donc il y a r\'eduction de cent vingt-huit dans la classe du c\^one. Les seize points singuliers de la surface donnent lieu \`a autant de lignes doubles dans le c\^one, ce qui effectue une r\'eduction de trente-deux; il y a encore une r\'eduction \`a effectuer de quatre-vingt-seize, qui doit avoir lieu \`a cause des lignes doubles ou des lignes de rebroussement du c\^one. En supposant qu'il y a $y$ de celles-ci et $x$ de celles-l\`a (outre les seize lignes doubles dont on a fait mention), il faut que l'on ait
\[
2x + 3y = 96,
\]
o\`u $x + y$ ne doit pas \^etre plus grand que trente-neuf; mais cela ne suffit pas pour d\'eterminer ces deux nombres.

Nous pouvons encore ajouter que les seize c\^ones qui touchent la surface aux points singuliers sont des c\^ones circonscrits, quatre \`a quatre, \`a quatre surfaces du second ordre, lesquelles sont aussi inscrites dans les seize plans singuliers situ\'ees quatre \`a quatre sur quatre surfaces du second ordre.

Je vais donner maintenant une id\'ee de la th\'eorie analytique. En repr\'esentant par
\[
x = 0,\ y = 0,\ z = 0,\ w = 0
\]
les \'equations des quatre faces du t\'etra\`edre, et par
\[
U = 0
\]
l'\'equation de la surface, $U$ sera une fonction homog\`ene du quatri\`eme degr\'e de $x, y, z, w$, laquelle, en faisant \'evanouir une quelconque des variables, doit se diviser en deux facteurs du second degr\'e, fonctions paires des trois autres variables; de plus, \`a cause de la condition par rapport aux points singuliers, $U$ ne peut pas contenir de terme $xyzw$. On doit donc avoir
\[
U = Ax^4 + By^4 + Cz^4 + Dw^4 + 2Fy^2 z^2 + 2Gx^2 z^2 + 2Hx^2 y^2 + 2Lx^2 w^2 + 2My^2 w^2 + 2Nz^2 w^2,
\]
o\`u il faut que les coefficients $A, B, C, F, G, H, L, M, N$ aient un syst\`eme inverse de la forme $0, 0, 0, 0, f, g, h, l, m, n$. Donc, en formant l'inverse de ce syst\`eme, on obtient, toute r\'eduction accomplie,
\[
U = mnfx^4 + nlgy^4 + lmhz^4 + fghw^4
\]
\[
\quad + (\ \ lf - mg - nh)(ly^2 z^2 + fx^2 w^2)
\]
\[
\quad + (-lf + mg - nh)(mz^2 x^2 + gy^2 w^2)
\]
\[
\quad + (-lf - mg + nh)(nx^2 y^2 + hz^2 w^2) = 0.
\]

En \'ecrivant
\[
\lambda = lf - mg - nh,\quad \mu = -lf + mg - nh,
\]
\[
\nu = -lf - mg + nh,\quad \omega = lf + mg + nh,
\]
\[
V = f^2 l^2 + m^2 g^2 + n^2 h^2 - 2mngh - 2nlhf - 2lmfg,
\]
l'\'equation de la courbe pourra s'\'ecrire sous les quatre formes
\[
(2mnfx^2 + n\nu y^2 + m\mu z^2 + f\lambda w^2)^2 = V(n^2 y^4 + m^2 z^4 + f^2 w^4 - 2mfy^2 w^2 - 2fn w^2 y^2 - 2nmy^2 z^2),
\]
\[
(n\nu x^2 + 2nlgy^2 + l^2\lambda z^2 + g\mu w^2)^2 = V(l^2 z^4 + g^2 w^4 + n^2 x^4 - 2gn w^2 x^2 - 2nl x^2 z^2 - 2lg z^2 w^2),
\]
\[
(m\mu x^2 + l^2\lambda y^2 + 2lmhz^2 + h\nu w^2)^2 = V(h^2 w^4 + m^2 x^4 + l^2 y^4 - 2ml x^2 y^2 - 2lh y^2 w^2 - 2hm w^2 x^2),
\]
\[
(f\lambda x^2 + g\mu y^2 + h\nu z^2 + 2fghw^2)^2 = V(f^2 x^4 + g^2 y^4 + h^2 z^4 - 2gh y^2 z^2 - 2hf z^2 x^2 - 2fg x^2 y^2);
\]
\footnote{En g\'en\'eral, quand deux syst\`emes de quantit\'es sont exprim\'es lin\'eairement les uns au moyen des autres, on dit que les deux syst\`emes de coefficients sont des syst\`emes inverses; on passe facilement de l\`a \`a l'id\'ee du syst\`eme inverse des coefficients d'une fonction du second ordre, et de tels syst\`emes se rencontrent si souvent, que l'on doit avoir un terme pour exprimer sans circonlocution cette relation.}

ce qui met en \'evidence les \'equations des sections de la surface par les quatre plans du t\'etra\`edre, et aussi celles des points singuliers, lesquelles sont, en effet,
\[
ny \pm mz \pm fw = 0,\quad lz \pm gw \pm nx = 0,
\]
\[
hw \pm mx \pm ly = 0,\quad fx \pm gy \pm hz = 0;
\]
et ces plans touchent la surface suivant les courbes d'intersection avec les surfaces
\[
2mnfx^2 + n\nu y^2 + m\mu z^2 + f\lambda w^2 = 0,\ \&c.,\ \&c.,
\]
ce qui d\'emontre le th\'eor\`eme \'enonc\'e par rapport aux seize courbes de contact des plans singuliers, et de l\`a celui pour les seize c\^ones tangents aux points singuliers.

Pour d\'eduire de l\`a la forme ordinaire de l'\'equation de la surface des ondes, \'ecrivons
\[
l = \alpha\beta\gamma(b\gamma - c\beta),\quad m = \alpha\beta\gamma(c\alpha - a\gamma),\quad n = \alpha\beta\gamma(a\beta - b\alpha),
\]
\[
f = ka\alpha(b\gamma - c\beta),\quad g = kb\beta(c\alpha - a\gamma),\quad h = kc\gamma(a\beta - b\alpha),
\]
\'equations qui suffisent pour d\'eterminer les rapports
\[
a : b : c : \alpha : \beta : \gamma : k
\]
au moyen de $l, m, n, f, g, h$. De cette mani\`ere, l'\'equation de la surface se r\'eduit \`a
\[
\alpha\beta\gamma(ax^2 + by^2 + cz^2)(\alpha x^2 + \beta y^2 + \gamma z^2)
\]
\[
\quad - ka\alpha(b\gamma + c\beta)x^2 w^2 - kb\beta(c\alpha + a\gamma)y^2 w^2 - kc\gamma(a\beta + b\alpha)z^2 w^2 + k^2 abc w^4 = 0,
\]
laquelle se r\'eduit \`a la surface des ondes en \'ecrivant
\[
x = \xi\sqrt{\alpha},\quad y = \eta\sqrt{\beta},\quad z = \zeta\sqrt{\gamma},
\]
$\xi, \eta, \zeta$ \'etant des coordonn\'ees rectangulaires, c'est-\`a-dire en faisant la transformation homographique du t\'etra\'edroide, de mani\`ere que l'un des plans du t\'etra\`edre passe \`a l'infini, et que les trois autres deviennent rectangulaires. De plus, en particularisant la transformation de mani\`ere que trois des coniques d'intersection se r\'eduisent \`a des cercles, cette surface rentre dans la surface des ondes. Il va sans dire que, dans le cas g\'en\'eral, plusieurs des points ou des plans dont nous avons parl\'e sont n\'ecessairement imaginaires; l'\'enum\'eration de tous les cas diff\'erents aurait \'et\'e d'une longueur effrayante. Il y aurait beaucoup \`a dire sur les cas particuliers o\`u quelques-unes des coniques se r\'eduisent \`a des paires de droites (r\'eelles ou imaginaires). Je me contente d'\'enoncer cette propri\'et\'e, tr\`es-facile \`a d\'emontrer, de la surface ordinaire des ondes: au cas o\`u $c = 0$, cette surface peut \^etre engendr\'ee par un cercle (ayant le centre de la surface pour centre, et dans un plan passant par l'axe des $z$), lequel se meut de mani\`ere \`a passer par la conique
\[
a^2 x^2 + c^2 y^2 = a^2 b^2.
\]
On trouve, dans le Cambridge and Dublin Mathematical Journal, t.~i.~[1846], p.~208, [38] la d\'emonstration d'une autre propri\'et\'e de la surface des ondes par rapport aux lignes de courbure des surfaces du second ordre.

\newpage

% ===========================================================================
% PDF page 328 -- Book page 306
% Paper [48]
% ===========================================================================
\section*{[48] Note sur les Fonctions de M. Sturm.}
\addcontentsline{toc}{section}{48. Note sur les Fonctions de M. Sturm}

\textit{[From the Journal de Math\'ematiques Pures et Appliqu\'ees (Liouville), tom.~xi.~(1846), pp.~297--299.]}

\bigskip

On sait que la suite des fonctions
\[
fx = (x - a)(x - b)(x - c)(x - d)\ldots,
\]
\[
f_1 x = \Sigma(x - b)(x - c)(x - d)\ldots,
\]
\[
f_2 x = \Sigma(a - b)^2(x - c)(x - d)\ldots,
\]
\[
f_3 x = \Sigma(a - b)^2 (b - c)^2 (c - a)^2 (x - d)\ldots,
\]
est de la plus grande utilit\'e dans la th\'eorie de la r\'esolution num\'erique des \'equations. En effet, on obtient tout de suite \`a leur moyen le nombre de racines r\'eelles comprises entre deux limites quelconques. Il \'etait donc int\'eressant de chercher la mani\`ere d'exprimer ces fonctions par les coefficients de $fx$.

Soit, pour cela, $m$ un nombre quelconque, pas plus grand que le degr\'e $n$ de cette fonction. En prenant $k$ pour $m$\textsuperscript{i\`eme} racine de la suite $a, b, c, \ldots,$ et mettant, pour abr\'eger,
\[
P = (-)^{\tfrac12 m(m-1)}(a - b)(a - c)\ldots(a - k)(b - c)\ldots(b - k)\ldots(j - k),
\]
cela donne
\[
f_m x : fx = \Sigma\,\frac{P^2}{(x - a)(x - b)\ldots(x - k)},
\]
dans laquelle expression
\[
P = \det\begin{pmatrix} 1 & a & \cdots & a^{m-1} \\ 1 & b & \cdots & b^{m-1} \\ \vdots & & & \vdots \\ 1 & k & \cdots & k^{m-1} \end{pmatrix}
\]
et, de plus,
\[
(x - a)(x - b)\ldots(x - k) = (-)^k \frac{m(m-1)\ldots(b - c)\ldots(b - k)\ldots(j - k)}{(x - a)}\cdots
\]
dans laquelle le coefficient de $x^{r-1}$ est \'egal \`a
\[
(-)^{\tfrac12 m(m-1)}\bigl[a^{m-1}(b - c)\ldots(b - k)\ldots(j - k) + \cdots\bigr].
\]
c'est-\`a-dire \`a
\[
(-)^{\tfrac12 m(m-1)}\,\det\begin{pmatrix} 1 & a & \cdots & a^{m-2} & a^{r-1} \\ 1 & b & \cdots & b^{m-2} & b^{r-1} \\ \vdots & & & & \vdots \\ 1 & k & \cdots & k^{m-2} & k^{r-1} \end{pmatrix}.
\]

Donc enfin le coefficient de $x^r$ dans $f_m x : fx$ est \'egal \`a
\[
\Sigma\,\det\begin{pmatrix} 1 & a & \cdots & a^{m-1} \\ 1 & b & \cdots & b^{m-1} \\ \vdots & & & \vdots \\ 1 & k & \cdots & k^{m-1} \end{pmatrix}\,\times\,\det\begin{pmatrix} 1 & a & \cdots & a^{m-2} & a^{r-1} \\ 1 & b & \cdots & b^{m-2} & b^{r-1} \\ \vdots & & & & \vdots \\ 1 & k & \cdots & k^{m-2} & k^{r-1} \end{pmatrix}
\]
ou, au moyen d'une propri\'et\'e connue des d\'eterminants, en repr\'esentant comme \`a l'ordinaire par $S_q$ la somme des $q$\textsuperscript{i\`emes} puissances de toutes les racines, ce coefficient devient \'egal \`a
\[
\det\begin{pmatrix} S_0 & S_1 & \cdots & S_{m-2} & S_{r-1} \\ S_1 & S_2 & \cdots & S_{m-1} & S_r \\ \vdots & & & & \vdots \\ S_{m-1} & S_m & \cdots & S_{2m-3} & S_{m+r-1} \end{pmatrix}.
\]

De l\`a, en mettant
\[
f_m x : fx = \Sigma\,\frac{T_q}{x - a}
\]
\[
= \Sigma a^{r}\cdots
\]
en multipliant par $fx$, et mettant
\[
Q_{m+s,r} = S_{m+r-1} - p_1 S_{m+r-2} + (-)^r p_r S_m,
\]
on obtient
\[
fx = x^n - p_1 x^{n-1}\ldots + (-)^n p_n,
\]
\[
f_m x = \Sigma_{r}^{n-m-1}\,\det\begin{pmatrix} S_0 & S_1 & \cdots & S_{m-2} & Q_{m,r} \\ S_1 & S_2 & \cdots & S_{m-1} & Q_{m+1,r} \\ \vdots & & & & \vdots \\ S_{m-1} & S_m & \cdots & S_{2m-3} & Q_{2m-1,r} \end{pmatrix}\,x^r
\]
o\`u $r$ peut ne s'\'etendre que depuis 0 jusqu'\`a $n - m$, puisque $f_m x$ est fonction enti\`ere. Au moyen des relations connues qui existent entre les quantit\'es $S_q$, on a
\[
Q_{m+s,r} = (-)^r p_{r-n} S_{m+s-r} \cdots + (-)^{m+r} p_n S_{r+m+s-n-1} \quad [r + m + s > n],
\]
\[
+ (-)^{r+m+s-n}\,p_{r+m+s-n-1}(r + m + s - 1)\,r + m + s - n,
\]
et de l\`a, en posant
\[
Q'_{m+s,r} = (-)^{r+m-1}\bigl\{p_{r+m} S_{s+1} \cdots + (-)^{m-1} p_n S_{r+m+s-1-n}\cdots\quad [r + m + s > n]
\]
\[
+ (-)^r p_{r+m+s-1}(r + m + s - n - 1)\bigr\},
\]
on peut, par les propri\'et\'es des d\'eterminants, r\'eduire $Q_{m+s,r}$ \`a $Q'_{m+s,r}$ dans l'expression de $f_m$. Nous avons donc exprim\'e cette fonction au moyen des coefficients $p_1, p_2, \ldots$ et des sommes $p$ ou $\Sigma$, $S_1, S_2, \ldots, S_{2m-1}$, lesquelles s'expriment elles-m\^emes par ces sommes (on aurait \`a calculer ces sommes une fois pour toutes jusqu'\`a $S_{2m-1}$); il serait donc facile de former des tables de ces fonctions, pour les \'equations d'un degr\'e quelconque, ce qui pourrait \`a peine s'effectuer d'aucune autre mani\`ere.

\newpage

% ===========================================================================
% PDF page 331 -- Book page 309
% Paper [49]
% ===========================================================================
\section*{[49] Sur quelques Formules du Calcul Int\'egral.}
\addcontentsline{toc}{section}{49. Sur quelques Formules du Calcul Int\'egral}

\textit{[From the Journal de Math\'ematiques Pures et Appliqu\'ees (Liouville), tom.~xii.~(1847), pp.~231--240.]}

\bigskip

Soit $x + y\sqrt{-1}$ ou $x + iy$ une quantit\'e imaginaire quelconque; faisons
\[
\rho = \sqrt{x^2 + y^2},\quad \theta = \arctan\frac{y}{x} \quad\cdots(1),
\]
$\rho$ \'etant une quantit\'e positive, et $\theta$ un arc compris entre les limites $\tfrac12\pi, -\tfrac12\pi$. Cela pos\'e, \'ecrivons
\[
(x + iy)^m = \rho^m e^{im\theta} \quad (x\ \text{positif}),
\]
\[
(x + iy)^m = \rho^m e^{im(\theta\pm\pi)} \quad (x\ \text{n\'egatif});
\]
(dans la seconde de ces formules, il faut prendre le signe sup\'erieur ou inf\'erieur, selon que $y$ est positif ou n\'egatif). Au cas de $x$ positif, la valeur du second membre sera ce que M.\ Cauchy a appel\'e valeur principale de $(x + iy)^m$. Au cas de $x$ n\'egatif, on peut aussi, \`a ce qu'il me semble, nommer cette valeur valeur principale. Cela para\^it contraire \`a la th\'eorie de M.\ Cauchy (\emph{Exercices de Math\'ematiques}, t.~i.~[1826] p.~2); mais la d\'emonstration que l'on y trouve de l'impossibilit\'e d'une valeur principale pour $x$ n\'egatif ne s'applique qu'au cas o\`u l'on suppose que le signe $\pm$ est toujours le m\^eme sans avoir \'egard au signe de $y$. Seulement, selon nos d\'efinitions, il importe de remarquer qu'il n'y a pas de valeur principale pour $x$ n\'egatif, au cas particulier o\`u $y = 0$; ou plut\^ot dans ce cas, et dans ce cas seulement, la valeur principale devient ind\'etermin\'ee.

Soit, en particulier, $x = 0$; les deux formules conduisent au m\^eme r\'esultat, savoir
\[
(iy)^m = (\pm y)^m e^{\pm im\pi/2} \quad\cdots(3).
\]
le signe comme auparavant. Car en consid\'erant $x$ comme infiniment petit positif, on obtient
\[
\theta = \pm\tfrac12\pi,
\]
et, en consid\'erant $x$ comme infiniment petit n\'egatif,
\[
\theta = \pm\tfrac12\pi,
\]
et de l\`a
\[
\theta \pm \pi = \mp\tfrac12\pi.
\]
Ainsi cette formule est toujours vraie, sans qu'il soit n\'ecessaire de consid\'erer $iy$ comme limite de $x + iy$, $x$ positif ou $x$ n\'egatif.

Remarquons encore que cette fonction $(iy)^m$ reste continue quand $y$ passe par z\'ero, ce qui a lieu aussi pour $(x + iy)^m$, $x$ positif, mais non pas pour $(x + iy)^m$, $x$ n\'egatif.

Les m\^emes remarques s'appliquent aux valeurs principales des logarithmes, lesquelles doivent se d\'efinir d'une mani\`ere analogue par les \'equations
\[
\log(x + iy) = \log\rho + i\theta\quad (x\ \text{positif}),
\]
\[
\log(x + iy) = \log\rho + i(\theta \pm \pi)\quad (x\ \text{n\'egatif})\quad\cdots(4).
\]
le signe ambigu, comme auparavant.

On d\'emontre sans difficult\'e que ces valeurs principales satisfont en tout cas aux \'equations
\[
(x + iy)^m (x' + iy')^m = \bigl[(x + iy)(x' + iy')\bigr]^m,
\]
\[
\log(x + iy) + \log(x' + iy') = \log\bigl[(x + iy)(x' + iy')\bigr]\quad\cdots(5);
\]
seulement ces \'equations deviennent ind\'etermin\'ees au cas o\`u, l'une des quantit\'es $x, x', xx' - yy'$ \'etant n\'egative, la quantit\'e correspondante $y, y', xy' + x'y$ s'\'evanouit.

Au moyen de cette d\'efinition de la valeur principale d'un logarithme, on obtient
\[
\int\,\frac{dx}{\alpha + \beta x} = \frac{A + B\theta}{\beta},\quad B = \log\Bigl(\frac{A + Bx}{\alpha}\Bigr) \quad\cdots(6),
\]
o\`u $\alpha, \beta$ sont r\'eels, et $A, B$ sont assujettis \`a la seule condition qu'au cas o\`u $\alpha, \beta$ seraient de signes contraires, la partie imaginaire de $\tfrac{A}{B}$ ne s'\'evanouisse pas. En effet, dans ce cas, l'int\'egrale et la valeur principale du logarithme deviennent toutes les deux ind\'etermin\'ees. Sans doute il y a une valeur que l'on peut appeler principale de l'int\'egrale, mais cette valeur n'est \'egale \`a aucun des logarithmes de $\tfrac{A + Bx}{\alpha + \beta x}$, et les notions des valeurs principales d'une int\'egrale et d'un logarithme n'ont pas de rapport ensemble. Ce r\'esultat s'accorde avec celui que j'ai trouv\'e dans mon M\'emoire ``Sur les fonctions doublement p\'eriodiques,'' [25].

Je passe \`a quelques autres applications de ces principes, qui ont rapport \`a la th\'eorie des fonctions $\Gamma$. Soit d'abord $r$ un nombre positif plus petit que l'unit\'e, et \'ecrivons
\[
U = \int_{-\infty}^{\infty}(ix)^{r-1} e^{ix}\,dx;
\]
on obtient
\[
U = e^{(r-1)\pi i/2}\int_0^\infty x^{r-1}\,e^{ix}\,dx + e^{-(r-1)\pi i/2}\int_0^\infty x^{r-1}\,e^{-ix}\,dx;
\]
ou enfin, puisque $x$ est positif dans ces deux int\'egrales,
\[
U = 2\cos(r - 1)\tfrac{\pi}{2}\,\Gamma r = 2\sin r\pi\,\Gamma r;
\]
savoir
\[
\sin r\pi\,\Gamma r = \tfrac12\int_{-\infty}^{\infty}(ix)^{r-1}\,e^{ix}\,dx.
\]
On obtiendrait de m\^eme
\[
0 = \tfrac12\int_{-\infty}^{\infty}(-ix)^{r-1}\,e^{ix}\,dx.
\]

Au premier coup d'\oe il, ces \'equations pourraient para\^itre en contradiction l'une avec l'autre: mais cela n'est pas ainsi, parce qu'il n'est pas vrai que $(-ix)^{r-1}$ soit \'egal \`a $(-i)^{r-1}(ix)^{r-1}$, $(-1)^{r-1}$ \'etant facteur invariable; mais au contraire $(-ix)^{r-1} = e^{\mp(r-1)\pi i}(ix)^{r-1}$, selon que $x$ est positif ou n\'egatif.

Dans la premi\`ere de ces \'equations, on peut remplacer $ix$ par $c + ix$, et dans la seconde, $-ix$ par $c - ix$, $c$ \'etant positif; mais cela n'est pas permis pour $c$ n\'egatif, \`a cause de la discontinuit\'e de valeur de $(c + ix)^{r-1}$ ou $(c - ix)^{r-1}$ dans ce cas, quand $x$ passe par z\'ero. En multipliant la premi\`ere \'equation par $e^{-c}$, et diff\'erentiant un nombre quelconque de fois, en admettant, pour les valeurs n\'egatives de $r$, l'\'equation
\[
\Gamma(r + 1) = r\,\Gamma r,
\]
la formule ne change pas de forme, et l'on obtient
\[
\sin r\pi\,\Gamma r\,e^{-c} = \tfrac12\int_{-\infty}^{\infty}(c + ix)^{r-1}\,e^{ix}\,dx \quad\cdots(7);
\]
de m\^eme, au moyen de la seconde \'equation,
\[
0 = \tfrac12\int_{-\infty}^{\infty}(c - ix)^{r-1}\,e^{ix}\,dx \quad\cdots(8),
\]
dans lesquelles formules $c$ est positif, et $r$ est un nombre quelconque entre $1, -\infty$, sans exclure la limite sup\'erieure; seulement, pour $r = 0$, le facteur $\sin r\pi\,\Gamma r$ se r\'eduit \`a $\pi$. Au cas o\`u $r$ est plus grand que z\'ero, on peut, si l'on veut, \'ecrire aussi $c = 0$. Dans tous les cas, on peut remplacer $\sin r\pi\,\Gamma r$ par $\tfrac{\pi}{\Gamma(1 - r)}$. Il importe de remarquer que ces m\^emes int\'egrales sont absolument inexprimables au cas o\`u $c$ est n\'egatif: en effet, en \'ecrivant $-c$ au lieu de $c$ ($c$ positif), on obtiendrait
\[
\int(-c + ix)^{r-1}\,e^{ix}\,dx = e^{(r-1)\pi i}\int(c - ix)^{r-1}\,e^{ix}\,dx + e^{-(r-1)\pi i}\int(c - ix)^{r-1}\,e^{ix}\,dx.
\]
Mais, par la seconde des \'equations dont il s'agit,
\[
0 = \tfrac12\int_{-\infty}^{\infty}(c - ix)^{r-1}\,e^{ix}\,dx + \tfrac12\int_{-\infty}^{\infty}(c - ix)^{r-1}\,e^{-ix}\,dx;
\]
donc
\[
\int(-c + ix)^{r-1}\,e^{ix}\,dx = -2i\sin\pi r\int(c - ix)^{r-1}\,e^{ix}\,dx.
\]
Or l'int\'egrale au second membre ne peut pas s'exprimer par les transcendantes connues, \`a moins que $r$ ne soit entier. \'Ecrivons encore, $c$ \'etant toujours positif,
\[
I_1 = \int(c + ix)^{r-1}\,e^{ix}\,dx,
\]
\[
I_2 = \int(c + ix)^{r-1}\,e^{-ix}\,dx,
\]
\[
I_3 = \int(c + ix)^{r-1}\,e^{ix}\,dx,
\]
\[
I_4 = \int(c - ix)^{r-1}\,e^{-ix}\,dx.
\]
Toutes les fonctions $\int(c \pm ix)^{r-1}\,e^{\pm ix}\,dx$ entre les limites $0, \infty$, ou $-\infty, 0$, ou $-\infty, \infty$, s'expriment facilement au moyen de $I_1, I_2, I_3, I_4$. Mais ces quatre fonctions ne sont pas connues; seulement, au moyen des \'equations qui viennent d'\^etre trouv\'ees, on obtient
\[
2\sin\pi r\,\Gamma r\,e^{-c} = I_1 + I_2,
\]
\[
0 = I_3 + I_4 \quad\cdots(10).
\]
On d\'eduit encore de ces m\^emes formules:
\[
\sin\pi r\,\Gamma r\,e^{-c}
\]
\[
= \int_{-\infty}^{\infty}(c + ix)^{r-1}\cos x\,dx = i\int_{-\infty}^{\infty}(c + ix)^{r-1}\sin x\,dx,
\]
\[
0 = \int(c - ix)^{r-1}\cos x\,dx = -i\int(c - ix)^{r-1}\sin x\,dx \quad\cdots(11).
\]

En supposant que $r - 1$ est entier n\'egatif, l'int\'egrale
\[
\int_{-\infty}^{\infty}(c^2 + x^2)^{r-1}\cos x\,dx
\]
se d\'ecompose facilement dans une suite d'int\'egrales de la forme
\[
\int_{-\infty}^{\infty}(c + ix)^{p-1}\cos x\,dx;
\]
et, en prenant la somme de celles-ci, on obtient la formule
\[
\int_{-\infty}^{\infty}(c^2 + x^2)^{r-1}\cos x\,dx = \frac{\sqrt{\pi}\,e^{-c}}{\Gamma(\tfrac32 - r)}\int_0^\infty(\xi + 2c)^{-r}\,e^{-\xi}\,d\xi \quad\cdots(12),
\]
laquelle est due \`a M.\ Catalan. Cependant, ni cette d\'emonstration ni celle de M.\ Catalan ne s'appliquent au cas o\`u $r$ n'est pas entier; la formule subsiste encore dans ce cas, comme M.\ Serret l'a d\'emontr\'e rigoureusement. En essayant de la v\'erifier, je suis tomb\'e sur cette autre formule:
\[
\int_{-\infty}^{\infty}(c^2 + x^2)^{r-1}\cos x\,dx
\]
\[
= \frac{\sqrt{\pi}\,e^{-c}\,(2c)^{r - 1/2}}{2\,\Gamma(1 - r)}\int_0^\infty\,\frac{(\sqrt{\xi + 2c} + \sqrt{\xi})^{1 - 2r} + (\sqrt{\xi + 2c} - \sqrt{\xi})^{1 - 2r}}{\sqrt{\xi + 2c}}\,e^{-\xi}\,d\xi \quad\cdots(13),
\]
ce qui suppose, comme auparavant, que $r - 1$ soit n\'egatif; en comparant les deux valeurs, on est conduit au r\'esultat singulier (en \'ecrivant $a$ au lieu de $2c$, et $\tfrac12(1 - p)$ pour $r$):
\[
\int_0^\infty\frac{(\sqrt{\xi + a} + \sqrt{\xi})^p + (\sqrt{\xi + a} - \sqrt{\xi})^p}{2\sqrt{\pi}}\,\xi^{(p-1)/2}\,e^{-\xi}\,d\xi = \frac{\Gamma\bigl(\tfrac{p+1}{2}\bigr) + a}{1 - p}\bigl(2\sqrt{a}\bigr)^p \quad\cdots(14),
\]
lequel peut se d\'emontrer sans difficult\'e, quand $p$ est entier positif impair, en d\'eveloppant les deux membres suivant les puissances de $\sqrt{a}/\sqrt{\xi}$, cela se fait au moyen de
\[
(\sqrt{1 + x} + 1)^p + (\sqrt{1 + x} - 1)^p = \sum 2^{p - 1 - r}\,\frac{\Gamma(p - r)}{\Gamma(p - 2r)\,\Gamma(r + 1)} \quad\cdots(15),
\]
o\`u $r$ s'\'etend depuis 0 jusqu'\`a $\tfrac12(p - 1)$. J'ai d\'eduit de cette formule (14) des formules assez remarquables qui se rapportent aux attractions, lesquelles para\^itront dans un num\'ero prochain du Cambridge and Dublin Mathematical Journal, [44]. On peut encore d\'emontrer cette formule singuli\`ere:
\[
\Gamma a\,\Gamma(r + a - 1) = \frac{1}{2\sin r\pi}\int_0^\infty (ix)^{a-1}(-ix)^{r-1}\,e^{ix}\,dx \quad\cdots(16);
\]
en effet, pour la v\'erifier, il suffit de r\'eduire l'int\'egrale, d'abord \`a
\[
\int(ix)^{a-1}(-ix)^{r-1}\,e^{ix}\,dx + \int(-ix)^{a-1}(ix)^{r-1}\,e^{-ix}\,dx,
\]
puis \`a
\[
\bigl(e^{(a-1)\pi i/2}\,e^{-(r-1)\pi i/2} + e^{-(a-1)\pi i/2}\,e^{(r-1)\pi i/2}\bigr)\int x^{a+r-2}\,e^{\pm ix}\,dx,
\]
dont les deux parties s'obtiennent au moyen d'une formule donn\'ee ci-dessus. Si, au lieu de $(-ix)^{r-1}$, l'on avait $(ix)^{r-1}$, l'int\'egrale se r\'eduirait \`a z\'ero; mais cela rentre dans une formule plus simple.

J'ajouterai encore ces deux formules-ci,
\[
2\pi c^{a-1}\,e^{-c} = \int_{-\infty}^{\infty}\,\frac{x^{a-1}}{c + ix}\,dx,
\]
\[
0 = \int_{-\infty}^{\infty}\,\frac{x^{a-1}\,e^{-ix}}{c + ix}\,dx,
\]
dont je supprime la d\'emonstration. En \'ecrivant dans la derni\`ere $-x$ au lieu de $x$, puis ajoutant, on a
\[
\pi c^{a-1}\,e^{-c} = \int_{-\infty}^{\infty}\,\frac{x^{a-1}}{c^2 + x^2}\,e^{ix}\,dx \quad\cdots(18);
\]
d'o\`u l'on d\'eduit tout de suite cette formule de M.\ Cauchy
\[
\tfrac12 \pi c^{a-1}\,e^{-c} = \int_{-\infty}^{\infty}\,\frac{x^{a-1}}{c^2 + x^2}\cos(\tfrac12 a\pi - x)\,dx \quad\cdots(19).
\]

La formule (7) peut \^etre consid\'er\'ee comme une d\'efinition de la fonction $\Gamma r$ au cas de $r$ n\'egatif; et, \`a ce point de vue, elle a, ce me semble, quelques avantages sur celle que M.\ Cauchy a donn\'ee au moyen des int\'egrales extraordinaires. Je passe \`a la d\'efinition, au moyen d'une int\'egrale d\'efinie, de la seconde int\'egrale eul\'erienne,
\[
\beta(m, n) = \frac{\Gamma m\,\Gamma n}{\Gamma(m + n)} \quad\cdots(20),
\]
quand $m$ ou $n$, ou tous les deux, sont n\'egatifs. Soit d'abord $m$ et $n$ tous les deux positifs, mais ni plus petits que l'unit\'e; on obtient au moyen de la formule (7) et de la d\'efinition ordinaire des fonctions $\Gamma$,
\[
\Gamma m\,\Gamma n = \pi^{-2}\int dx \int dy\,x^{m-1}(iy)^{n-1}\,e^{-(x+y)},
\]
et de l\`a, en mettant
\[
y = ax,\quad dy = x\,da,
\]
et int\'egrant par rapport \`a $x$,
\[
\beta(m, n) = \frac{1}{2i\sin n\pi}\int_{-\infty}^{\infty}(ia)^{n-1}\,(1 + a)^{-m-n}\,da,
\]
ou en \'ecrivant $k + ai$ au lieu de $ai$, $k$ \'etant positif,
\[
\beta(m, n) = \frac{1}{2i\sin n\pi}\int_{-\infty}^{\infty}\,\frac{(k + ai)^{n-1}}{(1 + k + ai)^{m+n}}\,da \quad\cdots(21),
\]
formule qui est vraie, m\^eme pour les valeurs n\'egatives de $n$. En effet, si cette formule est vraie pour une valeur quelconque particuli\`ere de $n$, elle sera aussi vraie pour la valeur $n + p$, $p$ \'etant entier positif quelconque; ce qui se d\'emontre au moyen des formules de r\'eduction: donc il ne s'agit que de la d\'emontrer au cas o\`u $n$ est positif et plus petit que l'unit\'e, et dans ce cas, puisque la fonction \`a int\'egrer est toujours finie, on peut \'ecrire sans crainte $k + ai$ au lieu de $ai$, ce qui la r\'eduit \`a la formule qui vient d'\^etre d\'emontr\'ee; donc cette formule (21) est la d\'efinition cherch\'ee, au cas o\`u $n$ est n\'egatif, ou positif et plus petit que l'unit\'e.

Si, de plus, $m$ est n\'egatif, ou positif et plus petit que l'unit\'e, $1 - m$ sera positif, et l'on d\'eduira
\[
\beta(m, n) = \frac{\Gamma m\,\Gamma n}{\Gamma(m + n)} = \frac{\Gamma(1 - m - n)\,\Gamma n\,\sin(m + n)\pi}{\Gamma(1 - m)\,\sin m\pi},
\]
c'est-\`a-dire
\[
\beta(m, n) = \frac{\sin(m + n)\pi}{\sin m\pi}\,\beta(1 - m - n, n);
\]
ou enfin, par l'\'equation (21),
\[
\beta(m, n) = \frac{\sin(m + n)\pi}{2i\sin m\pi\sin n\pi}\int_{-\infty}^{\infty}(k + ai)^{n-1}(1 + k + ai)^{-1+m}\,da,
\]
laquelle suppose seulement que $m + n$ soit n\'egatif, ou positif et plus petit que l'unit\'e. Elle pr\'esente une analogie assez frappante avec l'\'equation ordinaire
\[
\beta(m, n) = \int_0^1 a^{n-1}(1 - a)^{m-1}\,da
\]
qui correspond aux valeurs positives de $m, n$; de m\^eme que l'\'equation (21) est analogue \`a cette autre formule
\[
\beta(m, n) = \int_0^\infty\,\frac{a^{n-1}}{(1 + a)^{m+n}}\,da
\]
qui correspond aussi aux valeurs positives de $m$ et $n$.

On peut se proposer de v\'erifier l'\'equation ($m$ et $n$ positifs et plus petits que l'unit\'e)
\[
\Gamma m\,\Gamma n = \pi^{-2}\,\frac{1}{4\sin m\pi\sin n\pi}\int\!\!\int\,dx\,dy\,(ix)^{m-1}(iy)^{n-1}\,e^{i(x + y)},
\]
en transformant le second membre au moyen de $x = ay$. Pour cela, on distinguera quatre cas, selon que $x$ et $y$ sont tous les deux positifs ou n\'egatifs, ou l'un positif et l'autre n\'egatif. En mettant dans les deux premiers
\[
x = ay,\quad dx = y\,da,
\]
et en int\'egrant par rapport \`a $y$, on trouvera que les deux portions correspondantes de l'int\'egrale double se r\'euniront en
\[
-\Gamma(m + n)\cdot 2\cos(m + n)\pi\int_0^\infty\,\frac{a^{n-1}}{(1 + a)^{m+n}}\,da,
\]
c'est-\`a-dire \`a
\[
-2\cos(m + n)\pi\,\Gamma m\,\Gamma n.
\]
Pour les deux autres portions de l'int\'egrale double, en \'ecrivant
\[
x = ay,\quad dx = y\,da,
\]
on remarquera qu'il faut encore distinguer les deux cas $a < 1$ et $a > 1$. Les quatre int\'egrales ainsi obtenues se r\'euniront cependant dans les deux
\[
2\cos n\pi\,\Gamma(m + n)\int\,\frac{a^{n-1}\,da}{(1 + a)^{m+n}},\quad 2\cos m\pi\int\,\frac{a^{m-1}\,da}{(1 + a)^{m+n}},
\]
savoir, apr\`es quelques r\'eductions faciles, dans celles-ci,
\[
\frac{2\cos n\pi\sin m\pi}{\sin(m + n)\pi}\,\Gamma m\,\Gamma n,\quad \frac{2\cos m\pi\sin n\pi}{\sin(m + n)\pi}\,\Gamma m\,\Gamma n,
\]
lesquelles se r\'euniront en
\[
\cos(m - n)\pi\,\Gamma m\,\Gamma n.
\]
On obtient donc enfin l'\'equation identique
\[
4\sin m\pi\sin n\pi = 2\cos(m - n)\pi - 2\cos(m + n)\pi;
\]
ce qui suffit pour la v\'erification dont il s'agit.

\newpage

% ===========================================================================
% PDF page 339 -- Book page 317
% Paper [50]
% ===========================================================================
\section*{[50] Sur quelques Th\'eor\`emes de la G\'eom\'etrie de Position.}
\addcontentsline{toc}{section}{50. Sur quelques Th\'eor\`emes de la G\'eom\'etrie de Position}

\textit{[From the Journal f\"ur die reine und angewandte Mathematik (Crelle), tome xxxi.~(1846), pp.~213--227.]}

\bigskip

\subsection*{\S\,1.}

En prenant pour donn\'e un syst\`eme quelconque de points et de droites, on peut mener par deux points donn\'es des nouvelles droites, ou trouver des points nouveaux, savoir les points d'intersection de deux des droites donn\'ees; et ainsi de suite. On obtient de cette mani\`ere un nouveau syst\`eme de points et de droites, qui peut avoir la propri\'et\'e que plusieurs des points sont situ\'es dans une m\^eme droite, ou que plusieurs des droites passent par le m\^eme point; ce qui donne lieu \`a autant de th\'eor\`emes de g\'eom\'etrie de position. On a d\'ej\`a \'etudi\'e la th\'eorie de plusieurs de ces syst\`emes; par exemple de celui de quatre points; de six points, situ\'es deux \`a deux sur trois droites qui se rencontrent dans un m\^eme point; de six points trois \`a trois sur deux droites, ou plus g\'en\'eralement, de six points sur une conique (ce dernier cas, celui de l'hexagramme mystique de Pascal, n'est pas encore \'epuis\'e; nous y reviendrons dans la suite), et m\^eme de quelques syst\`emes dans l'espace. Cependant il existe des syst\`emes plus g\'en\'eraux que ceux qui ont \'et\'e examin\'es, et dont les propri\'et\'es peuvent \^etre aper\c{c}ues d'une mani\`ere presque intuitive, et qui, \`a ce que je crois, sont nouveaux. Commen\c{c}ons par le cas le plus simple. Imaginons un nombre $n$ de points situ\'es d'une mani\`ere quelconque dans l'espace, et que nous d\'esignerons par $1, 2, 3, \ldots, n$. Qu'on fasse passer par toutes les combinaisons de deux points des droites, et par toutes les combinaisons de trois points des plans; puis coupons ces droites et ces plans par un plan quelconque, les droites selon des points, et les plans selon des droites. Soit $\alpha\beta$ le point qui correspond \`a la droite men\'ee par les deux points $\alpha, \beta$; soit de m\^eme $\beta\gamma$ le point qui correspond \`a celle men\'ee par les points $\beta, \gamma$, et ainsi de suite. Soit de plus $\alpha\beta\gamma$ la droite qui correspond au plan passant par les trois points $\alpha, \beta, \gamma$, etc. Il est clair que les trois points $\alpha\beta, \alpha\gamma, \beta\gamma$ seront situ\'es dans la droite $\alpha\beta\gamma$. Donc en repr\'esentant par $N_2, N_3,\ldots$ les nombres des combinaisons de $n$ lettres prises deux \`a deux, trois \`a trois etc.\ \`a la fois, on a le th\'eor\`eme suivant:

\textbf{Th\'eor\`eme I.} \emph{On peut former un syst\`eme de $N_2$ points situ\'es trois \`a trois sur $N_3$ droites: savoir en repr\'esentant les points par $12, 13, 23$, etc.\ et les droites par $123$, etc., les points $12, 13, 23$ seront situ\'es sur la droite $123$, et ainsi de suite.}

Pour $n = 3$, ou $n = 4$, cela est tout simple; on aura trois points sur une droite, ou six points trois \`a trois sur quatre droites; il n'entre dans aucune propri\'et\'e g\'eom\'etrique. Pour $n = 5$ on a dix points
\[
12, 13, 14, 15, 23, 24, 25, 34, 35, 45
\]
et les droites
\[
123, 124, 125, 134, 135, 145, 234, 235, 245, 345.
\]

Les points $12, 13, 14, 23, 24, 34$ sont les angles d'un quadrilat\`ere quelconque\footnote{Il faut avoir \'egard toujours \`a la diff\'erence entre quadrilat\`ere et quadrangle; chaque quadrilat\`ere a quatre c\^ot\'es et six angles, chaque quadrangle a quatre angles et six c\^ot\'es.}; le point 15 est tout \`a fait arbitraire, le point 25 est situ\'e sur la droite passant par les points 12 et 15, mais sa position sur cette droite est arbitraire. On d\'eterminera depuis les points $35, 45$; 35 comme point d'intersection des droites passant par 13 et 15 et par 23 et 25, c'est-\`a-dire des droites 135 et 235, et de m\^eme 45 comme point d'intersection des lignes 145 et 245. Les points 35 et 45 auront la propri\'et\'e g\'eom\'etrique d'\^etre en ligne droite avec 34, ou bien tous les trois seront dans une m\^eme droite 345.

\'Etudions de plus pr\`es la figure que nous venons de former. En prenant les cinq num\'eros dans un ordre d\'etermin\'e, par exemple dans l'ordre naturel $1, 2, 3, 4, 5$, les cinq points $12, 23, 34, 45, 51$ pourront \^etre consid\'er\'es comme formant un pentagone que nous repr\'esenterons par la notation $(12345)$. Les c\^ot\'es de ce pentagone sont \'evidemment $123, 234, 345, 451, 512$. De m\^eme les points $13, 35, 52, 24, 41$ peuvent \^etre consid\'er\'es comme formant le pentagone $(13524)$ dont les c\^ot\'es sont $135, 352, 524, 241, 413$. Ce pentagone est circonscrit au premier, car ses c\^ot\'es passent \'evidemment par les angles $15, 23, 45, 12, 34$ du premier; mais il est de m\^eme inscrit \`a celui-ci, car ses angles sont situ\'es respectivement dans les c\^ot\'es $123, 345, 512, 234, 451$ de ce m\^eme pentagone. Donc les pentagones
\[
(12345),\quad (13524)
\]
sont \`a la fois circonscrits et inscrits l'un \`a l'autre, donc:

\textbf{Th\'eor\`eme II.} \emph{La figure compos\'ee de dix points, trois \`a trois dans dix droites, peut \^etre consid\'er\'ee (m\^eme de six mani\`eres diff\'erentes) sous la forme de deux pentagones, inscrits et circonscrits l'un \`a l'autre.}

Ou encore

\textbf{Th\'eor\`eme III.} \emph{\'Etant donn\'e un pentagone quelconque, on peut toujours trouver un autre pentagone qui y est \`a la fois circonscrit et inscrit. Ce second pentagone peut satisfaire \`a une seule condition donn\'ee quelconque.}

Si par exemple le second pentagone doit avoir un de ses angles sur un point donn\'e d'un c\^ot\'e du premier, la construction se d\'eduit tout de suite de ce qui pr\'ec\`ede.

Ces paires correspondantes de pentagones forment une figure connue. On en trouve la construction dans une note de M.~[J.~T.] Graves dans le \emph{Philosophical Magazine} [vol.~xv.~1839], mais la m\^eme figure est encore mieux connue sous un autre point de vue. En effet, consid\'erons le point 12, et les droites 123, 124, 125 qui passent par ce point; puis les triangles dont les angles sont 13, 14, 15 et 23, 24, 25. Les c\^ot\'es de ces m\^emes triangles sont 134, 135, 145 et 234, 235, 245, et les c\^ot\'es correspondants se rencontrent dans les points 34, 35, 45 qui sont en ligne droite. Donc le th\'eor\`eme sur les pentagones est le suivant:

``Si les angles de deux triangles sont situ\'es deux \`a deux dans trois droites qui se rencontrent dans un point, leurs c\^ot\'es homologues se coupent dans trois points en ligne droite.''

Remarquons aussi que ce th\'eor\`eme particulier (en n'empruntant rien des trois dimensions de l'espace) reproduit le th\'eor\`eme g\'en\'eral relatif au nombre $n$. Il n'y a pour cela qu'\`a consid\'erer $n$ droites passant par le m\^eme point, et qui peuvent \^etre d\'esign\'ees par $1, 2, 3, \ldots, n$. En choisissant d'abord les points 12, 13, tout triangle dont les trois angles sont situ\'es dans les droites $1, 2, 3$, pendant que deux de ses c\^ot\'es passent par 12, 13, a la propri\'et\'e que le troisi\`eme c\^ot\'e passe par un point d\'etermin\'e 23 situ\'e dans la droite passant par 12, 13. En prenant arbitrairement le point 14, on obtient avec les droites $1, 3, 4$ ou $1, 2, 4$ les nouveaux points 34, 24 qui sont en ligne droite avec 23, et ainsi de suite.

Passons au cas $n = 6$. Il existe ici quinze points situ\'es trois \`a trois sur vingt droites, ou bien vingt droites qui se coupent quatre \`a quatre en quinze points. Il n'y a point ici des syst\`emes d'hexagones, mais il existe un syst\`eme de neuf points qui est assez remarquable. Divisons d'une mani\`ere quelconque les num\'eros 1, 2, 3, 4, 5, 6 en deux suites par trois, par exemple en 1, 3, 5 et 2, 4, 6, et consid\'erons les neuf points
\[
12,\ 14,\ 16,
\]
\[
32,\ 34,\ 36,
\]
\[
52,\ 54,\ 56.
\]
Les droites qui passent par 12 et 32, 14 et 34, 16 et 36, savoir 132, 134, 136, se rencontrent dans le m\^eme point 13. De m\^eme les droites qui passent par 32 et 52, 34 et 54, 36 et 56 se rencontrent dans 35, et les droites qui passent par 12 et 52, 14 et 54, 16 et 56 se rencontrent dans 15. Les points 13, 15 et 35 sont sur la m\^eme droite 135. En consid\'erant les points 12, 14, 16 comme formant un triangle, et de m\^eme les points 32, 34, 36 et 52, 54, 56, cela revient \`a dire que les droites men\'ees par les angles homologues des triangles prises deux \`a deux, se rencontrent trois \`a trois dans trois points situ\'es dans la m\^eme droite. Ou bien, ce que l'on savait d\'ej\`a par le th\'eor\`eme 3: les c\^ot\'es homologues des triangles se rencontrent trois \`a trois dans trois points situ\'es en ligne droite. En effet, les c\^ot\'es des triangles sont 124, 126, 146 pour la premi\`ere, et 324, 326, 346 et 524, 526, 546 pour les deux autres. Les trois premiers c\^ot\'es se rencontrent dans 24, les autres dans 26 et 46, et ces trois points sont dans la droite 246. Maintenant tout cela arrive \'egalement en combinant les colonnes verticales, ou en consid\'erant les neuf points comme formant les trois autres triangles dont les angles sont 12, 32, 52; 14, 34, 54; 16, 36, 56. Cela donne lieu au th\'eor\`eme suivant:

\textbf{Th\'eor\`eme IV.} \emph{Le syst\`eme de quinze points, situ\'es trois \`a trois sur vingt droites, contient (et cela m\^eme de dix mani\`eres diff\'erentes) un syst\`eme de neuf points qui ont la propri\'et\'e de former de deux mani\`eres diff\'erentes trois triangles, tels, que les droites qui passent par leurs angles homologues, prises deux \`a deux, se rencontrent dans trois points qui sont en ligne droite, tandis que les c\^ot\'es homologues des triangles se coupent trois \`a trois en trois autres points qui sont aussi en ligne droite. Dans la seconde mani\`ere de former les triangles, ces deux syst\`emes de trois points en ligne droite sont seulement \'echang\'es.}

Il ne reste qu'\`a savoir combien il y en a d'arbitraires dans le syst\`eme de quinze points situ\'es trois \`a trois sur vingt droites. En supposant le syst\`eme form\'e pour le nombre cinq, on peut prendre arbitrairement 16 et 26 sur la droite 126 qui est d\'etermin\'ee par les points 12 et 16. Donc 12, 13, 14, 15 et 16 sont arbitraires et 23, 24, 25, 26 sont arbitrairement situ\'es sur des droites donn\'ees. L'existence des droites 345, 346, 356, 456 constitue autant de th\'eor\`emes g\'eom\'etriques; c'est-\`a-dire, chacune de ces droites est d\'etermin\'ee par trois points.

En essayant d'approfondir la th\'eorie de six points sur la m\^eme conique, on rencontrera un syst\`eme de neuf points, tel que ceux que nous venons d'examiner; mais il est moins g\'en\'eral. Il existe des relations entre les points qui n'ont pas lieu dans le syst\`eme g\'en\'eral. Je renvoie cette discussion \`a une section s\'epar\'ee de ce m\'emoire, et je passe au cas de $n = 7$.

Pour ce cas on a tout de suite le th\'eor\`eme suivant:

\textbf{Th\'eor\`eme V.} \emph{Le syst\`eme de vingt et un points situ\'es trois \`a trois sur trente-cinq droites, peut \^etre consid\'er\'e (m\^eme de cent vingt mani\`eres diff\'erentes) comme compos\'e de trois heptagones, le premier circonscrit au second, le second au troisi\`eme et le troisi\`eme au premier. Les heptagones par exemple peuvent \^etre $(1234567), (1357246), (1526374)$.}

Dans ce syst\`eme 12, 13, 14, 15, 16, 17 sont arbitraires, et 23, 24, 25, 26, 27 le sont sur des droites donn\'ees; les droites 345, 346, 347, 356, 357, 367, 456, 457, 467, 567 sont d\'etermin\'ees chacune par trois points. Dans le cas g\'en\'eral 12, 13, $\ldots$, 1$n$ sont arbitraires, et $23 \ldots 2n$ le sont sur des droites donn\'ees. Il existe $\tfrac16(n - 2)(n - 3)(n - 4)$ droites dont chacune est d\'etermin\'ee par trois points. Un th\'eor\`eme analogue \`a celui-ci a lieu quand $n$ est un nombre premier; savoir le suivant:

\textbf{Th\'eor\`eme VI.} \emph{Le syst\`eme de $N_2$ points, situ\'es trois \`a trois sur $N_3$ droites, peut \^etre consid\'er\'e (m\^eme de $\frac{1 \cdot 2 \ldots (n-2)}{n-1}$ mani\`eres) comme compos\'e de $\tfrac12(n - 1)$ $n$-gones, le premier circonscrit au second, le second au troisi\`eme, etc., et le dernier au premier.}

Je ne connais pas d'autres cas o\`u l'id\'ee des nombres premiers se pr\'esente dans la g\'eom\'etrie. Il sera peut-\^etre possible de trouver des th\'eor\`emes analogues \`a III, IV, V, pour toutes les formes du nombre $n$, mais je n'ai pas encore examin\'e cela.

Le th\'eor\`eme g\'en\'eral I, peut \^etre consid\'er\'e comme l'expression d'un fait analytique, qui doit \'egalement avoir lieu en consid\'erant quatre coordonn\'ees au lieu de trois. Ici une interpr\'etation g\'eom\'etrique a lieu, qui s'applique aux points dans l'espace. On peut en effet, sans recourir \`a aucune notion m\'etaphysique \`a l'\'egard de la possibilit\'e de l'espace \`a quatre dimensions, raisonner comme suit (tout cela pourra aussi \^etre traduit facilement en langue purement analytique): En supposant quatre dimensions de l'espace, il faudra consid\'erer des lignes d\'etermin\'ees par deux points, des demi-plans d\'etermin\'es par trois points, et des plans d\'etermin\'es par quatre points; (deux plans se coupent alors selon un demi-plan, etc.). L'espace ordinaire doit \^etre consid\'er\'e comme plan, et il coupera un plan selon un plan ordinaire, un demi-plan selon une ligne ordinaire, et une ligne selon un point ordinaire. Tout cela pos\'e: en consid\'erant un nombre $n$ de points, et les combinant deux \`a deux, trois \`a trois, et quatre \`a quatre par des lignes, des demi-plans et des plans, puis coupant le syst\`eme par l'espace consid\'er\'e comme plan, on obtient le th\'eor\`eme suivant de g\'eom\'etrie \`a trois dimensions:

\textbf{Th\'eor\`eme VII.} \emph{On peut former un syst\`eme de $N_2$ points, situ\'es trois \`a trois dans $N_3$ droites qui elles-m\^emes sont situ\'ees quatre \`a quatre dans $N_4$ plans. En repr\'esentant les points par 12, 13, etc., les points situ\'es dans la m\^eme droite sont 12, 13, 23; et les droites \'etant repr\'esent\'ees par 123 etc.\ comme auparavant, les droites 123, 124, 134, 234 sont situ\'ees dans le m\^eme plan 1234.}

En coupant cette figure par un plan, on obtient le th\'eor\`eme suivant de g\'eom\'etrie plane:

\textbf{Th\'eor\`eme VIII.} \emph{On peut former un syst\`eme de $N_3$ points situ\'es quatre \`a quatre dans $N_4$ droites. Alors 123, 124, 134, 234 sont dans la m\^eme droite d\'esign\'ee par 1234.}

De m\^eme, en consid\'erant un espace \`a $p + 2$ dimensions, on obtient la proposition suivante, encore plus g\'en\'erale:

\textbf{Th\'eor\`eme IX.} \emph{On peut former dans l'espace un syst\`eme de $N_p$ points, qui passent $p + 1$ \`a $p + 1$ par $N_{p+1}$ droites, situ\'ees $p + 2$ \`a $p + 2$ dans $N_{p+2}$ plans, ou bien pour la g\'eom\'etrie plane, un syst\`eme de $N_p$ points, situ\'es $p + 1$ \`a $p + 1$ dans $N_{p+1}$ droites.}

Des th\'eor\`emes analogues \`a IV et V seraient probablement tr\`es nombreux et tr\`es compliqu\'es.

Les r\'eciproques polaires auront \'evidemment lieu pour tous ces th\'eor\`emes; on pourrait aussi les d\'emontrer directement d'une mani\`ere analogue.

\subsection*{\S\,II. Sur le Th\'eor\`eme de Pascal.}

En consid\'erant six points sur la m\^eme conique, et les prenant dans un ordre d\'etermin\'e, pour en former un hexagone, on sait que les c\^ot\'es oppos\'es se rencontrent dans trois points situ\'es en ligne droite. En prenant les points dans un ordre quelconque, on en peut former soixante hexagones, \`a chacun desquels correspond une droite; il s'agit maintenant de trouver les relations entre ces droites.

M.~Steiner a prouv\'e dans son ouvrage \emph{Systematische Entwickelungen u.s.w.}~[1832], que ces soixante droites passent trois \`a trois par vingt points, et il ajoute que ces vingt points sont situ\'es quatre \`a quatre sur quinze droites. La premi\`ere partie de ce th\'eor\`eme peut \^etre d\'emontr\'ee assez facilement, comme nous le verrons: mais pour la seconde partie, je n'ai pas r\'eussi \`a trouver les combinaisons de quatre points qui doivent \^etre situ\'es en ligne droite, et il me para\^it m\^eme qu'il est impossible de les trouver\footnote{Je ne sais pas s'il existe une d\'emonstration de la seconde partie du th\'eor\`eme; je n'ai pu la trouver nulle part. Au cas que cette partie du th\'eor\`eme n'\'etait pas correcte, il para\^it que l'on devra peut-\^etre lui substituer la proposition suivante: ``Les vingt points d\'eterminent deux \`a deux dix lignes qui passent trois \`a trois par dix points.'' On verra dans ce qui suit, de quelle mani\`ere il faudrait combiner ces points. [Voir 55.]}.

Cherchons les combinaisons des droites qui doivent passer trois \`a trois par le m\^eme point.

Soient 1, 2, 3, 4, 5, 6 les six points situ\'es sur la m\^eme conique. Consid\'erons d'abord l'hexagone 123456 que l'on obtient en prenant les points dans un ordre d\'etermin\'e. Suivant le th\'eor\`eme de Pascal les trois points
\[
12.45,\quad 23.56,\quad 34.61
\]
(o\`u, 12.45 d\'esigne le point d'intersection des lignes passant par les points 1, 2 et $\{4, 5\}$) sont situ\'es en ligne droite. Consid\'erons les six hexagones
\[
12\ 3\ 4\ 5\ 6
\]
\[
14\ 3\ 6\ 5\ 2
\]
\[
16\ 3\ 2\ 5\ 4
\]
\[
14\ 3\ 2\ 5\ 6
\]
\[
12\ 3\ 6\ 5\ 4
\]
\[
16\ 3\ 4\ 5\ 2
\]
qu'on tire du premier en permutant les nombres 2, 4, 6 correspondants aux sommets altern\'es de l'hexagone. Pour les trois premiers on fait les permutations cycliques de ces nombres (savoir 246, 462, 624), pour les trois autres on fait d'abord une inversion 426, puis les permutations cycliques (426, 264, 642). En \'ecrivant les combinaisons des points qui doivent \^etre situ\'es en ligne droite, on a
\begin{center}
\begin{tabular}{l l l}
$12.45$ & $23.56$ & $34.61$\\
$36.12$ & $56.14$ & $52.34$\\
$45.36$ & $14.23$ & $16.52$\\[2pt]
$14.25$ & $43.56$ & $32.61$\\
$36.14$ & $56.12$ & $54.42$\\
$25.36$ & $12.34$ & $61.54$\\
\end{tabular}
\end{center}
Suivant cette table les points sur la m\^eme horizontale sont en ligne droite.

On remarquera d'abord que les trois premi\`eres droites passent par les angles des triangles dont les c\^ot\'es sont 36, 45, 12 et 14, 23, 56. Les c\^ot\'es homologues de ces triangles se rencontrent en 36.14, 45.23, 12.56 qui sont en ligne droite, c'est-\`a-dire, par un th\'eor\`eme d\'ej\`a cit\'e: les trois lignes passent par un m\^eme point. On aurait \'et\'e conduit au m\^eme r\'esultat en observant que les trois premi\`eres droites passent par les triangles dont les c\^ot\'es sont 14, 23, 56 et 52, 16, 34 ou enfin 52, 16, 34 et 36, 45, 12. De m\^eme les trois derni\`eres droites passent par le m\^eme point. Donc il a \'et\'e d\'emontr\'e ce qui suit:

\textbf{Th\'eor\`eme X.} \emph{En consid\'erant les trois hexagones qu'on obtient en permutant cycliquement les angles altern\'es du premier, les trois droites qui y correspondent se rencontrent dans un m\^eme point. Les soixante lignes passent donc trois \`a trois par vingt points.}

Ajoutons qu'aux trois hexagones de ce th\'eor\`eme correspondent d'une mani\`ere particuli\`ere trois autres hexagones, ou que les vingt points doivent se combiner deux \`a deux d'une mani\`ere particuli\`ere.

Mais on se formera une id\'ee plus claire du syst\`eme en remarquant que les neuf droites
\[
36,\ 45,\ 12
\]
\[
14,\ 23,\ 56
\]
\[
25,\ 61,\ 34
\]
ont entre elles une relation qui est polaire r\'eciproque de celle entre les neuf points du th\'eor\`eme IV. Pour faciliter cette comparaison, je prendrai d'abord le th\'eor\`eme analogue pour les tangentes d'une conique.

\textbf{Th\'eor\`eme XI.} \emph{Soient $1, 3, 5$ et $2, 4, 6$ des tangentes \`a une m\^eme conique et $12$, etc.\ les points d'intersection de ces droites: les neuf points
\[
36,\ 45,\ 12
\]
\[
14,\ 23,\ 56
\]
\[
25,\ 61,\ 34
\]
peuvent \^etre d\'etermin\'es au moyen de six points de l'espace $A, B, C, \alpha, \beta, \gamma$, de mani\`ere que $A\alpha$, etc.\ repr\'esente le point d'intersection de la droite passant par $A, \alpha$ avec le plan de la figure. Les points sont correspondants entre eux de cette mani\`ere:
\[
A\alpha,\ A\beta,\ A\gamma
\]
\[
B\alpha,\ B\beta,\ B\gamma
\]
\[
C\alpha,\ C\beta,\ C\gamma
\]
seulement les points 36, 23, 34 etc.\ sont en ligne droite, ce qui n'aurait pas lieu pour les points $A\alpha, B\beta, C\gamma$, si la position de $A, B, C, \alpha, \beta, \gamma$ \'etait arbitraire. On est donc conduit \`a ce probl\`eme:}

Trouver six points $A, B, C, \alpha, \beta, \gamma$ dans l'espace, tels, qu'en repr\'esentant par $A\alpha$, etc.\ l'intersection de la droite men\'ee par $A\alpha$ avec un plan donn\'e, les combinaisons des points
\[
(A\alpha,\ B\beta,\ C\gamma)
\]
\[
(A\beta,\ B\gamma,\ C\alpha)
\]
\[
(A\gamma,\ B\alpha,\ C\beta)
\]
\[
(A\alpha,\ B\gamma,\ C\beta)
\]
\[
(A\beta,\ B\alpha,\ C\gamma)
\]
\[
(A\gamma,\ B\beta,\ C\alpha)
\]
soient en ligne droite.

Pour le th\'eor\`eme de Pascal, cela donne:

\textbf{Th\'eor\`eme XII.} \emph{Soient $1, 3, 5$ et $2, 4, 6$ des points d'une conique, les neuf lignes
\[
36,\ 45,\ 12
\]
\[
14,\ 23,\ 56
\]
\[
25,\ 61,\ 34
\]
peuvent \^etre consid\'er\'ees comme les projections des lignes
\[
A\alpha,\ A\beta,\ A\gamma
\]
\[
B\alpha,\ B\beta,\ B\gamma
\]
\[
C\alpha,\ C\beta,\ C\gamma
\]
sur le plan de la figure, o\`u $A, B, C, \alpha, \beta, \gamma$ sont six plans, dont la relation reste encore \`a d\'eterminer.}

En effectuant la solution du probl\`eme que j'ai indiqu\'ee on aurait, \`a ce qu'il me semble, un point de vue tout \`a fait nouveau d'envisager les coniques.

Je vais ajouter encore quelques r\'eflexions sur la mani\`ere de chercher les relations qui existent entre les vingt points. En \'ecrivant seulement les angles altern\'es des hexagones, on a cette table:
\[
1.2.3
\]
\[
1.2.4
\]
\[
1.2.5
\]
\[
1.2.6
\]
\[
1.3.4
\]
\[
1.3.5
\]
\[
1.3.6
\]
\[
1.4.5
\]
\[
1.4.6
\]
\[
1.5.6
\]
\`A chaque symbole correspondent six hexagones, qui, \`a ce que nous avons vu, se partagent en deux paires de trois hexagones, et \`a chaque combinaison de trois, il correspond un point. Il y a donc deux points qui correspondent au symbole 1.3.5, deux qui correspondent au symbole 1.3.6, deux au symbole 1.5.6 etc. En repr\'esentant donc par 35, 36, 56, les droites passant par ces paires de points, il me para\^it probable que ces droites aient ensemble les relations du th\'eor\`eme I, (savoir que 35, 36, 56 se rencontrent dans un point etc.), ce qui donnerait lieu au th\'eor\`eme hypoth\'etique que j'ai \'enonc\'e dans une note. Voil\`a, \`a ce que je puis apercevoir, la seule mani\`ere sym\'etrique de combiner les droites. Mais au moins les symboles
\[
1.3.5
\]
\[
1.3.6
\]
\[
1.5.6
\]
ont entre eux des rapports singuliers. En effet, \'ecrivons pour chacun les neuf points du th\'eor\`eme XII, on a ce tableau:
\begin{center}
\begin{tabular}{l l l}
$36$, & $45$, & $12$\\
$14$, & $23$, & $56$\\
$25$, & $61$, & $34$\\[3pt]
$35$, & $64$, & $12$\\
$14$, & $23$, & $56$\\
$26$, & $15$, & $34$\\[3pt]
$35$, & $46$, & $12$\\
$14$, & $25$, & $36$\\
$26$, & $13$, & $54$\\
\end{tabular}
\end{center}
qui ne contient que quatorze points. Cela m\'erite des recherches ult\'erieures.

\medskip
\noindent\textbf{D\'emonstration analytique du th\'eor\`eme de Pascal, et de la premi\`ere partie de celui de M.~Steiner. Formules relatives au m\^eme sujet.}
\medskip

Soient $P = 0,\ Q = 0,\ R = 0$ les \'equations des lignes 12, 34, 56. On d\'emontrera assez facilement que les \'equations des lignes 45, 61, 23 peuvent \^etre repr\'esent\'ees par
\[
P + \nu Q + \mu R = 0,
\]
\[
\nu P + Q + \lambda R = 0,
\]
\[
\mu P + \lambda Q + R = 0.
\]

En effet les six points 1, 2, 3, 4, 5, 6 seront situ\'es dans la conique
\[
P^2 + Q^2 + R^2 + \frac{\lambda + \tfrac{1}{\lambda}}{1}\,QR + \frac{\mu + \tfrac{1}{\mu}}{1}\,PR + \frac{\nu + \tfrac{1}{\nu}}{1}\,PQ = 0;
\]
car en faisant dans cette \'equation $P = 0$, l'\'equation se r\'eduit \`a
\[
(\mu Q + \lambda R)(\lambda Q + R) = 0;
\]
c'est-\`a-dire, la conique contient les points d\'etermin\'es par
\[
(P = 0,\ \mu P + Q + \lambda R = 0),
\]
\[
(P = 0,\ \mu P + \lambda Q + R = 0),
\]
ou bien les points 1, 2; et de m\^eme elle contient les autres points 3, 4, 5, 6. Les fonctions $P, Q, R$ sont cens\'ees contenir chacune deux constantes arbitraires; donc on a neuf constantes arbitraires dans ce syst\`eme, qui par cons\'equent est tout-\`a-fait g\'en\'eral. On peut former le syst\`eme suivant d'\'equations:
\begin{align*}
12.\quad & P = 0,\\
13.\quad & \lambda\mu P + Q + \lambda R = 0,\\
14.\quad & \lambda P + \mu Q + \lambda\mu R = 0,\\
15.\quad & P + \nu Q + \nu\lambda R = 0,\\
16.\quad & \nu P + Q + \lambda R = 0,\\
23.\quad & \lambda\mu P + \lambda Q + R = 0,\\
24.\quad & P + \mu\lambda Q + \mu R = 0,\\
25.\quad & \lambda P + \nu\lambda Q + \nu R = 0,\\
26.\quad & \nu\lambda P + \lambda Q + R = 0,\\
34.\quad & Q = 0,\\
35.\quad & \mu P + \mu\nu Q + R = 0,\\
36.\quad & \mu\nu P + \mu Q + \nu R = 0,\\
45.\quad & P + \nu Q + \mu R = 0,\\
46.\quad & \nu P + Q + \nu\mu R = 0,\\
56.\quad & R = 0.
\end{align*}

\'Ecrivons les \'equations des lignes comprises dans la table de neuf points ci-dessus donn\'es. On a d'abord
\[
\mu\nu P + \mu Q + \nu R = 0,\quad P + \nu Q + \mu R = 0,\quad P = 0,
\]
\[
\lambda P + \mu Q + \lambda\mu R = 0,\quad \mu P + \lambda Q + R = 0,\quad R = 0,
\]
\[
\lambda P + \nu\lambda Q + \nu R = 0,\quad \nu P + Q + \lambda R = 0,\quad Q = 0.
\]
En combinant la seconde et la troisi\`eme colonne verticale du tableau, on obtient pour les trois points d'intersection des c\^ot\'es oppos\'es de l'hexagone 123456, les \'equations
\[
(P = 0,\ \mu Q + \nu R = 0),
\]
\[
(R = 0,\ \lambda P + \mu Q = 0),
\]
\[
(Q = 0,\ \lambda P + \nu R = 0),
\]
qui appartiennent \`a trois points situ\'es sur la droite
\[
\lambda P + \mu Q + \nu R = 0,
\]
ce qui suffit pour d\'emontrer le th\'eor\`eme de Pascal.

On obtient de m\^eme, en combinant les autres paires de colonnes verticales, deux syst\`emes de trois points, respectivement situ\'es dans les droites
\[
\frac{\lambda}{\mu} P + \frac{\mu}{\nu} Q + \frac{\nu}{\lambda} R = 0,
\]
\[
\lambda\mu\nu P + \mu Q + \nu R = 0,
\]
lesquelles, avec la droite qu'on vient de trouver,
\[
\lambda P + \mu Q + \nu R = 0,
\]
se rencontrent \'evidemment dans un m\^eme point, d\'etermin\'e par les deux \'equations
\[
\lambda P + \mu Q + \nu R = 0\quad\text{et}\quad \frac{P}{\lambda} = \frac{Q}{\mu} = \frac{R}{\nu}.
\]

Voil\`a une d\'emonstration de la premi\`ere partie du th\'eor\`eme de M.~Steiner. Les \'equations que nous venons de trouver appartiennent au point d'intersection des trois droites qui correspondent au premier des trois hexagones du symbole 1.3.5. Pour trouver l'autre point correspondant de la m\^eme mani\`ere \`a ce symbole, il faut combiner les colonnes horizontales, ce qui donne pour les coordonn\'ees de ce point:
\[
(\lambda - \mu\nu) P = (\mu - \nu\lambda) Q = (\nu - \lambda\mu) R.
\]

En cherchant de m\^eme les expressions des points qui correspondent aux symboles 1.3.6 et 1.5.6, on obtient des r\'esultats moins \'el\'egants, mais qui valent peut-\^etre la peine d'\^etre \'enonc\'es ici.

Je forme cette table compl\`ete:

\bigskip
\noindent\textbf{Pour 1.3.5}
\begin{center}
Syst\`emes de trois lignes, qui se rencontrent dans un point.
\end{center}
\[
\text{I.}\quad\begin{cases} \lambda P + \mu Q + \nu R = 0,\\ \dfrac{P}{\lambda} = \dfrac{Q}{\mu} = \dfrac{R}{\nu}.\end{cases}
\]
\[
\text{II.}\quad \bigl(\lambda P + \mu Q + \nu R\bigr) + \lambda\mu\nu\Bigl(\frac{P}{\lambda} + \frac{Q}{\mu} + \frac{R}{\nu}\Bigr) = 0;
\]

\bigskip
\noindent\textbf{Pour 1.3.6}
\[
\text{I.}\quad\begin{cases} (\lambda - \mu\nu) P = (\nu - \lambda\mu) R,\\ (\nu - \lambda\mu) R = (\mu - \nu\lambda) Q,\\ (\mu - \nu\lambda) Q = (\lambda - \mu\nu) P; \end{cases}
\]
\[
\text{II.}\quad\begin{cases} \mu P + \lambda Q + \lambda\mu\nu R = 0,\\ \nu\lambda P + \mu\nu Q + R = 0,\\ (\mu P + \lambda Q + \lambda\mu\nu R) + (\nu\lambda P + \mu\nu Q + R) = 0; \end{cases}
\]

\bigskip
\noindent\textbf{Pour 1.5.6}
\[
\text{I.}\quad\begin{cases} (\mu - \nu\lambda) P + (1 - \lambda\mu\nu) R = 0,\\ (1 - \nu\lambda\mu) R + (\lambda - \mu\nu) Q = 0,\\ (\lambda - \mu\nu) Q - (\mu - \nu\lambda) P = 0; \end{cases}
\]
\[
\text{II.}\quad\begin{cases} (\mu\nu - \lambda\mu\nu^2 - \lambda + \lambda\mu^2) P + (\mu - \nu\lambda) Q + (\mu^2\nu - \lambda\mu\nu^2 + \mu - \lambda\mu) R = 0,\\ (\lambda - \lambda\mu^2 + \mu\nu - \lambda\mu^2) P + (\mu - \lambda\mu - \nu) Q + (\nu - \lambda\mu\nu^2) R = 0,\\ (\lambda\nu^2 - \lambda\mu^2\nu^2 - \mu\nu + \lambda\mu^2) P + (\nu\lambda - \mu\nu^2) Q + (\lambda\mu - \mu^2\nu) R = 0. \end{cases}
\]

\bigskip
\noindent\textbf{Note sur le th\'eor\`eme de M.~Brianchon.}
\medskip

On peut donner une d\'emonstration semblable de ce th\'eor\`eme, en prenant pour les \'equations des six tangentes celles-ci:
\[
1.\ P = 0,
\]
\[
3.\ Q = 0,
\]
\[
5.\ R = 0,
\]
\[
4.\ \alpha P + \beta Q + \gamma R = 0,
\]
\[
6.\ \alpha' P + \beta' Q + \gamma' R = 0,
\]
\[
2.\ \alpha'' P + \beta'' Q + \gamma'' R = 0,
\]
et en cherchant la relation entre les coefficients qui est n\'ecessaire pour que ces six \'equations appartiennent aux tangentes d'une m\^eme conique. On obtient facilement
\[
(\alpha\gamma' - \alpha'\gamma)(\beta'\alpha'' - \beta''\alpha')(\gamma''\beta - \gamma\beta'') = (\alpha'\gamma'' - \alpha''\gamma')(\beta''\alpha - \beta\alpha'')(\gamma\beta' - \gamma'\beta),
\]
ce qui exprime aussi la condition pour que les trois diagonales se rencontrent dans un m\^eme point.

\selectlanguage{english}

\end{document}
